\documentclass[11pt,reqno]{amsart}
\usepackage[T1]{fontenc}
\usepackage[utf8]{inputenc}
\usepackage{lmodern}
\usepackage[letterpaper,margin=1.05in]{geometry}
\usepackage{amsmath,amssymb,amsfonts,mathtools,mathrsfs}
\usepackage{microtype}
\usepackage{enumitem,booktabs,longtable,array}
\usepackage[dvipsnames]{xcolor}
\usepackage{hyperref}
\hypersetup{colorlinks=true,linkcolor=MidnightBlue,citecolor=MidnightBlue,urlcolor=MidnightBlue,
  pdftitle={Projective division and the homogeneous Arveson--Douglas conjecture},
  pdfsubject={Research manuscript on the homogeneous Arveson--Douglas conjecture}}
\numberwithin{equation}{section}
\newtheorem{theorem}{Theorem}[section]
\newtheorem{lemma}[theorem]{Lemma}
\newtheorem{proposition}[theorem]{Proposition}
\newtheorem{corollary}[theorem]{Corollary}
% The five principal results retain the labels of the announcement.
\newtheorem{thmA}{Theorem}\renewcommand{\thethmA}{A}
\newtheorem{thmB}{Theorem}\renewcommand{\thethmB}{B}
\newtheorem{thmC}{Theorem}\renewcommand{\thethmC}{C}
\newtheorem{thmD}{Theorem}\renewcommand{\thethmD}{D}
\newtheorem{thmE}{Theorem}\renewcommand{\thethmE}{E}

\theoremstyle{definition}
\newtheorem{definition}[theorem]{Definition}
\newtheorem{example}[theorem]{Example}
\theoremstyle{remark}
\newtheorem{remark}[theorem]{Remark}
\newcommand{\C}{\mathbb C}
\newcommand{\R}{\mathbb R}
\newcommand{\N}{\mathbb N}
\newcommand{\Z}{\mathbb Z}
\newcommand{\Pj}{\mathbb P}
\newcommand{\cO}{\mathcal O}
\newcommand{\cI}{\mathcal I}
\newcommand{\cK}{\mathcal K}
\newcommand{\cS}{\mathcal S}
\newcommand{\FS}{\mathrm{FS}}
\newcommand{\DA}{\mathrm{DA}}
\newcommand{\dbar}{\bar\partial}
\newcommand{\eps}{\varepsilon}
\newcommand{\id}{\mathrm{id}}
\newcommand{\norm}[1]{\left\lVert#1\right\rVert}
\newcommand{\abs}[1]{\left\lvert#1\right\rvert}
\DeclareMathOperator{\im}{im}
\DeclareMathOperator{\rank}{rank}
\DeclareMathOperator{\supp}{supp}
\DeclareMathOperator{\Supp}{Supp}
\DeclareMathOperator{\Tr}{Tr}
\DeclareMathOperator{\Hom}{Hom}
\DeclareMathOperator{\End}{End}
\DeclareMathOperator{\Gr}{Gr}
\DeclareMathOperator{\Quot}{Quot}
\DeclareMathOperator{\Hilb}{Hilb}
\DeclareMathOperator{\Span}{span}
\DeclareMathOperator{\dist}{dist}
\DeclareMathOperator{\diam}{diam}
\DeclareMathOperator{\Proj}{Proj}
\DeclareMathOperator{\Spec}{Spec}
\DeclareMathOperator{\coker}{coker}
\allowdisplaybreaks[2]
\setlist[enumerate]{label=\textup{(\roman*)},leftmargin=2.4em}
\setlist[itemize]{leftmargin=2em}
\emergencystretch=1.4em
\title[Projective division and the homogeneous Arveson--Douglas conjecture]{Projective division and the homogeneous\\Arveson--Douglas conjecture}
\date{24 November 2026}
\author{Lijia Ding}
\address{School of Mathematics and Statistics, Zhengzhou University, 450001, Zhengzhou, P.R.China}
\email{ljding@zzu.edu.cn}
\author{Xiang Fang}
\address{Department of Applied Mathematics, National Yang Ming Chiao Tung University, Hsinchu, Taiwan}
\email{xfang@nycu.edu.tw}
\author{Xiang Tang}
\address{Department of Mathematics, Washington University, St. Louis, MO, 63130, U.S.A.}
\email{xtang@math.wustl.edu}
\thanks{Research manuscript for independent mathematical assessment. The proof concerns the standard homogeneous conjecture. }
\subjclass[2020]{Primary 47A13; Secondary 32W05, 13D02, 47B10}
\keywords{Arveson--Douglas conjecture, Drury--Arveson space, homogeneous ideals,
  coherent sheaves, polynomial division, essential normality, Schatten ideals}
\begin{document}
\begin{abstract}
We develop a proof of the homogeneous Arveson--Douglas conjecture for
submodules of the finite-multiplicity Drury--Arveson shift. For a fixed
homogeneous submodule, the positive square of the column of coordinate
components discarded by the quotient projection is bounded by $C/(n+1)$
on degree $n$. If the algebraic quotient has support dimension $\delta>0$,
its coordinate commutators belong to the weak Schatten class
$\cS_{\delta,\infty}$ and hence to $\cS_p$ for every $p>\delta$.
The argument has three stages. We first realize the scalar graded problem
on projective space with its exact Hilbert norm and construct complete
presentations stable under normalized limits. We then prove a uniform
holomorphic lifting estimate using quantitative division, sampled Hermitian
metrics, and a finite $\bar\partial$ correction. Finally we return the
estimate to the original shift and transfer it to finite coefficient
multiplicity before counting singular values. The projective constructions use the full ideal sheaves, including their
nonreduced structure; saturation removes only finite-length information
at the cone vertex.
We explain the relation to the classical division, resolution, and
Dolbeault methods, and identify the simultaneous degree-uniform estimates
on which the argument depends.
\end{abstract}
\maketitle
\tableofcontents

% BEGIN sections/introduction.tex
\section{Introduction}
\label{sec:introduction}

The essential normality problem asks whether polynomial relations preserve,
up to compact operators, the commutation relations between the coordinates
of a holomorphic Hilbert space and their adjoints.  Homogeneity makes this
question quantitative.  Each commutator preserves total degree, and the
degree-$n$ space is finite dimensional.  One must control the norm of the
commutator on that space and the number of vectors on which it acts.  The
first is an analytic problem; the second is measured by the Hilbert function
of the algebraic quotient.


% BEGIN sections/history.tex
\subsection{Historical context and the role of division}
\label{sec:history}

Arveson's development of the multivariable shift made the space now called
the Drury--Arveson space a basic model for commuting row contractions.
His 1998 paper established the Schatten estimates for the full shift and
the associated Toeplitz extension
\cite[Proposition~5.3 and Theorem~5.7]{Arveson1998}.
The passage to polynomially defined subspaces introduced a further
question: how much of this analytic structure survives compression?
For submodules of a finite-multiplicity shift generated by vector
monomials, Arveson proved that the projection commutators lie in
$\cS_p$ for $p>2m$, and that the quotient coordinate commutators lie in
$\cS_p$ for $p>m$.  His Conjecture~A asks for compactness of the
projection commutators for arbitrary homogeneous vector generators
\cite[Theorem~2.1, Corollary~2.2 and Section~5]{Arveson2005}.
Thus finite coefficient multiplicity and the off-diagonal projection
operators belong to the original formulation.

Douglas connected essential normality with the construction of
$K$-homology classes and proposed that the quotient's dimension should
determine its summability threshold.  In \emph{A new kind of index
theorem}, Conjecture~9 states this refinement for the quotient of the
ball Bergman module by the closure of a polynomial ideal: the proposed
threshold is $p>\dim(Z\cap\mathbb B^m)$, where $Z$ is the ideal's zero
variety \cite{DouglasIndex}.  The corresponding dimension refinement in
the homogeneous Drury--Arveson setting is the one considered here.
Guo and Kai Wang established important cases, including homogeneous
principal submodules in arbitrary ambient dimension and scalar homogeneous
quotients in two and three variables \cite{GuoWang}.
These results supplied both analytic estimates and concrete connections
between quotient modules and algebraic geometry.

Shalit made norm-controlled polynomial division an explicit organizing
principle \cite{ShalitDivision}.  His stable division property asks for
generators $f_1,\ldots,f_s$ and decompositions
$h=\sum a_\nu f_\nu$ satisfying
$\sum\|a_\nu f_\nu\|\le C\|h\|$ in the relevant norm.
The proof of his Theorem~4.1 differentiates such a decomposition,
projects away terms still in the module, and obtains an off-diagonal
bound of order $(n+1)^{-1/2}$.  The quotient Hilbert function then gives
the dimension-sensitive Schatten conclusion.  This division-and-projection
mechanism is a direct antecedent of the analytic reduction below.
There is also a specific antecedent for the scalar-to-vector passage:
Section~5 of \cite{ShalitDivision} encodes a vector polynomial by a
scalar polynomial linear in auxiliary variables.  Its equation~(5.1)
records the degree-dependent factor relating the two coordinate shifts.
The transfer in Section~\ref{sec:completion} retains that factor and
performs the final singular-value count in the original quotient.

Eschmeier's operator-inequality criterion supplies a particularly close
antecedent \cite{EschmeierCriterion}.  Its projected-derivative form is
recorded explicitly in \cite[Section~6, equation~(6.1)]{ShalitDivision}:
for homogeneous generators $f_\nu$ of common degree $D$, the estimate
\[
 \left\|P_{\overline I^\perp}\sum_\nu a_\nu\partial_{z_j}f_\nu\right\|
 \le C\sqrt n\left\|\sum_\nu a_\nu f_\nu\right\|,
 \qquad \deg a_\nu=n-D,
\]
uniformly in the homogeneous coefficients and the degree, implies
essential normality of the quotient with Schatten threshold
$p>\dim(R/I)$.  This criterion identifies the projected derivative as
the analytic quantity that must be controlled.  The lifting construction
below is designed to establish precisely this type of estimate for
arbitrary homogeneous relations in the original Drury--Arveson norm.

The geometric result of Engli\v{s} and Eschmeier is particularly relevant
\cite{EnglisEschmeier}.  For a homogeneous variety $V$ smooth away from
the origin, they considered the submodule of functions vanishing on
$V\cap\mathbb B^m$; algebraically this is the radical-ideal setting.
Their argument uses restriction to the variety and the generalized
Toeplitz calculus of Boutet de Monvel and Guillemin.  It covers the
Drury--Arveson space within a larger scale of holomorphic Sobolev spaces.
For $\delta=\dim V$, they obtain quotient commutators in $\cS_p$ for
$p>\delta$.  They also explicitly record
$[S_i,S_j^*]\in\cS_{\delta,\infty}$ immediately after the homogeneous
main proof in Section~5.  The weak Schatten endpoint in this geometric
setting is therefore already part of their result.  Their treatment
also allows smooth components of different dimensions meeting only at
the origin.  This contribution provides an essential reference point
for both the hypotheses and the quantitative conclusion of the present
manuscript.

Douglas, Tang, and Yu brought the analysis of the $\bar\partial$ equation
on singular complex spaces into this index-theoretic problem
\cite{DouglasTangYu}. Under their assumptions of $M\le m-2$ defining equations, a Jacobian of
rank $M$ along the boundary, and sphere transversality, the zero set in the ball is a complete
intersection with smooth strongly pseudoconvex boundary and at most
isolated interior singularities. Their Dolbeault--Neumann analysis and
relative $K$-homology identify a boundary fundamental class; the bounded
restriction and extension theorem (extending Beatrous' smooth transverse
case \cite{Beatrous}) connects that geometric class to the quotient of
the ball Bergman module. Their Theorems 1.2, 1.4, and 1.5 establish the
corresponding relative class, bounded extension, and essential normality.
These results do not assume homogeneity. The analytic location of the
present argument is different: the $\bar\partial$ equations below are
solved on the smooth ambient projective space, in bundles presenting an
ideal sheaf, rather than on the regular locus of its zero variety.

The resolution method in operator theory is another important antecedent.
Douglas, Jabbari, Tang, and Yu \cite{DouglasJabbariTangYu} realize monomial
ideals through finite exact complexes of essentially normal Hilbert modules
and use additivity to compute the associated quotient $K$-homology class.
Their method also applies in the Drury--Arveson setting. The complexes in
this article are initially sheaf complexes; algebraic exactness is not
being identified with exactness after Hilbert completion. A quantitative
lifting theorem is needed to obtain operator estimates from the sheaf presentation.
Wang's asymptotic stable division method \cite{YiWang} similarly uses local
norm-controlled decompositions to study Bergman submodules. The local-to-global
strategy is thus well established in essential normality.

\subsection{Division and Dolbeault methods in complex geometry}
The analytic literature provides two further lines of direct antecedents.
First, quantitative standard-basis division of convergent series allows
both coefficients and equations to vary. Galligo's parameter-dependent
division theorem \cite[Theorems~1.2.7--1.2.8]{Galligo} gives estimates for
quotients and remainders on suitable polydiscs, and relates zero remainders
to flatness. Hauser and M\"uller \cite[Theorems~5.1--5.2]{HauserMueller}
obtain continuous linear scissions for analytic module maps, with estimates
on weighted coefficient spaces. Their method compares a monomial splitting
with a perturbed reconstruction operator. Shalit's weighted leading-term
argument \cite[Section~3]{ShalitDivision} has the same relevance in the
graded Hilbert-module problem. Theorem C below uses this classical
mechanism and proves the homogeneous-family version needed here. In
particular, choosing the homogenizing variable to have unit weight permits
a degree-independent affine estimate. Theorem C is a division statement
for image sheaves, not an extension theorem for sections of a quotient
sheaf on a subvariety.

Second, homological $\bar\partial$ methods for division go back at least
to H\"ormander \cite{HormanderGenerators} and Kelleher--Taylor
\cite[Section~II]{KelleherTaylor}. Their use of a Koszul complex to
organize successive corrections is a precursor of the descent below.
Skoda's ideal-division theorem \cite{SkodaDivision} instead uses direct
weighted $L^2$ estimates; his later work on surjective vector bundle maps
\cite{SkodaBundles} treats holomorphic lifting with positivity assumptions.
Twisted $L^2$ division and division in generically exact complexes are
studied by Varolin \cite{Varolin} and Ji \cite{Ji}. These results make the
choice of coefficient metrics and the control of curvature central issues.
Section~\ref{ana:metrics} explains why the precise unweighted datum norm
required here cannot simply be replaced by a singular generator weight.

Particularly close homological precedents for projective lifting are Andersson
\cite[Section~2]{AnderssonMembership} and Andersson--Wulcan
\cite[Section~6, Proposition~6.1]{AnderssonWulcan}. They combine projective
resolutions with Dolbeault equations and positive twists to solve polynomial
membership problems. In the latter work, residue currents encode the
membership conditions and the successive equations of the complex. Here
uniform local division of holomorphic overlap data supplies smooth
Dolbeault forms instead; the subsequent descent is a quantitative version
of the familiar homological correction. Neither the use of a resolution
nor the formal sequence of $\bar\partial$ equations is in itself a new
principle.

The additional assertion of the present article is the compatibility of
all the estimates with the original homogeneous Hilbert norms. Complete
presentations in uniformly regular degrees retain the polynomial image
spaces at normalized limits, giving one uniform division estimate for
all centers and scales. Sampled coefficient metrics control local division,
differentiated generators, and curvature simultaneously. At the scale
$\varepsilon=R_*/\sqrt n$, positive twisting and finite descent produce
a global holomorphic lift without changing the scalar norm of the datum.
The derivative estimate returns to the original Hilbert norm before any
singular values are counted. We prove the particular statements B--E
below rather than assume an analytic admissibility or stable-division
condition on the given ideal. The article does not claim novelty for a
classical construction in isolation; it is the quantitative combination
of the constructions that supports the asserted homogeneous theorem.

% END sections/history.tex


\subsection{The statement}
Let $m,r\ge1$, let $R=\C[z_1,\ldots,z_m]$, and let $H_m^2$ be the Drury--Arveson space,
normalized by
\[
 \langle z^\alpha,z^\beta\rangle=0\quad(\alpha\ne\beta),
 \qquad \|z^\alpha\|^2=\frac{\alpha!}{|\alpha|!}.
\]
Fix a finite-dimensional Hermitian coefficient space $\C^r$, placed in
degree zero.  Let $M\subset R\otimes\C^r$ be a homogeneous submodule and set
\[
 \mathcal M=\overline M\subset H_m^2\otimes\C^r,\qquad
 Q=\mathcal M^\perp,\qquad Q_n=H_{m,n}^2\otimes\C^r\ominus M_n.
\]
Write $Z_j=M_{z_j}\otimes I_r$ and
\[
 S_j=P_QZ_j|_Q,\qquad V_j=Z_j|_{\mathcal M},\qquad
 B_{j,n}=P_{M_{n+1}}Z_j|_{Q_n}.
\]
The positive operator
\begin{equation}\label{eq:intro-defect}
 A_n=\sum_{j=1}^m B_{j,n}^*B_{j,n}
\end{equation}
measures the total squared component removed by the quotient projection.

\begin{thmA}[Homogeneous essential normality in finite multiplicity]\label{thm:main}
For every homogeneous submodule $M\subset R\otimes\C^r$ there is a
constant $C_M$ such that
\begin{equation}\label{eq:main-degree}
 0\le A_n\le\frac{C_M}{n+1}I_{Q_n}\qquad(n\ge0).
\end{equation}
Consequently
\begin{equation}\label{eq:main-commutator}
 \|[S_i^*,S_j]|_{Q_n}\|\le\frac{C_M'}{n+1}
 \qquad(1\le i,j\le m).
\end{equation}
If $T=(R\otimes\C^r)/M$ is infinite dimensional and
$\delta=\dim\Supp T$, then
\[
 [S_i^*,S_j]\in\cS_{\delta,\infty}\subset\bigcap_{p>\delta}\cS_p.
\]
The column $B=(P_{\mathcal M}Z_j|_Q)_{j=1}^m$ belongs to
$\cS_{2\delta,\infty}$, as do the individual off-diagonal operators
$[P_{\mathcal M},Z_j]$ after extension by zero.  The restriction
commutators satisfy
\[
 [V_i^*,V_j]\in\cS_{m,\infty}\subset\bigcap_{p>m}\cS_p.
\]
If $T$ is finite dimensional, the quotient commutators and off-diagonal
operators have finite rank; the restriction assertion still holds.
\end{thmA}

Here $\cS_{p,\infty}$ consists of compact operators whose singular values
satisfy $s_k=O(k^{-1/p})$.  The constant in
\eqref{eq:main-degree} depends on the fixed module and its original
Hermitian geometry.  The theorem includes every closed submodule invariant
under the standard circle action, since its polynomial part is finitely
generated and dense; this is proved in Lemma~\ref{op:graded}.

For a scalar ideal $I$, the number $\delta$ is the complex dimension of
$V(I)$, whether or not $I$ is radical.  Thus the theorem includes arbitrary
singular homogeneous varieties and homogeneous nonreduced structures.
The proof concerns the standard grading.  It makes no assertion for a
general nonhomogeneous ideal or for a weighted grading with unequal positive
weights.

\subsection{The three stages of the proof}
The proof separates an operator-norm estimate from the algebraic count of
its multiplicities. Section~\ref{sec:operators} gives the exact compression
identity: it suffices to estimate the column $B_n$ by $O(n^{-1/2})$ in the
original Hilbert space. We first prove that estimate for scalar ideals.

\medskip\noindent\emph{Step 1: From the graded Hilbert module to projective space.}
Section~\ref{sec:projective} identifies a homogeneous polynomial of degree
$n$ with a section of $L^n$ over $W=\Pj^{m-1}$, $L=\cO_W(1)$. Saturation
changes only finitely many blocks. The two Hilbert norms are related by
\[
 \|F_n\|_{\DA}^2=\binom{n+m-1}{m-1}\|F_n\|_{\FS}^2,
\]
so the projections agree within each degree. Section~\ref{sec:algebra},
culminating in Theorem~\ref{thm:B}, constructs complete presentations of the
associated ideal sheaf and controls their exactness and Hilbert functions
under the rescalings used in the analysis. This is the algebraic part of
the uniformity argument.

\medskip\noindent\emph{Step 2: Uniform division and $\bar\partial$ analysis.}
Theorem~\ref{div:uniform} establishes a uniform two-ball estimate for
compact coefficient families. It is applied at each stage of the
presentation. Theorem~\ref{thm:D} constructs coefficient metrics that
control differentiation and curvature and make the differentials
contractive. Sections~\ref{ana:local-and-global}--\ref{ana:global-section}
then prove Theorem~\ref{geo:lifting}, the bounded global lifting theorem.
The desired output, for $F\in H^0(W,\mathscr I\otimes L^{n+1})$, is
\[
 \Phi a=F,\qquad
 \|a\|_{H_{0,\varepsilon}\otimes L^{n+1}}\le C\|F\|_{L^{n+1}},
 \qquad\varepsilon=R_*/\sqrt n,
\]
with $|(\partial_{z_j}\Phi)v|_{L^{-1}}\le
C\varepsilon^{-1}|v|_{H_{0,\varepsilon}}$ in the very same metric.
The $\bar\partial$ equations are on the smooth presenting bundles.
One differentiated cutoff costs $O(\varepsilon^{-1})$ and one positive-twist
solution gains $O(\varepsilon)$; a finite descent cancels those powers.

\medskip\noindent\emph{Step 3: Back to the Drury--Arveson shift.}
Section~\ref{sec:completion} differentiates $F=\Phi a$ and projects
perpendicular to the ideal, eliminating $\Phi\,\partial a$.
The exact adjacent-degree norm ratio yields the desired $O(n^{-1/2})$
scalar off-diagonal bound. A square-zero auxiliary-variable encoding
then gives the vector estimate before the Hilbert function is used.
Thus the exponent is the dimension of the original support, not the
ambient dimension of the auxiliary scalar problem.

Section~\ref{sec:extensions} records the consequences for other standard
ball kernels and analytic ideals generated by homogeneous functions, and
an example showing why the weak Schatten endpoint cannot generally be
made strong. The appendix specifies the classical theorems and records
the dependencies of the quantitative argument. The notation A--E agrees
with the research announcement.

We use $\N=\{0,1,2,\ldots\}$.  All dimensions of algebraic supports are complex dimensions.  All sheaf
maps are algebraic on projective space and are also regarded as holomorphic
maps after analytification.  Constants may depend on the fixed ideal,
presentation, charts and chosen radii; they never depend on the large
degree or on the section being lifted.  A zero coefficient stage is omitted
and all subsequent stages are taken to be zero.  The one-variable scalar
case is elementary and is treated separately.

% END sections/introduction.tex


% BEGIN sections/operators.tex
\section{Graded quotients and coordinate defects}
\label{sec:operators}

We first separate the operator estimate from the algebraic count.  This
also fixes the original Hilbert-space norms in which every later
comparison will be made.

\subsection{The grading and its polynomial part}
For a multi-index $\alpha$, put $|\alpha|=\sum_j\alpha_j$ and
$\alpha!=\prod_j\alpha_j!$.  The vectors
\[
 e_\alpha=\sqrt{\frac{|\alpha|!}{\alpha!}}z^\alpha
\]
form an orthonormal basis of $H_m^2$.  If $|\alpha|=n$, then
\begin{equation}\label{eq:op-shift}
 Z_je_\alpha=\sqrt{\frac{\alpha_j+1}{n+1}}e_{\alpha+e_j},
 \qquad
 Z_j^*z^\alpha=\frac{\alpha_j}{n}z^{\alpha-e_j}\quad(n>0),
\end{equation}
with the second expression interpreted as zero if $\alpha_j=0$.
In particular $\|Z_j\|\le1$.  Tensoring with $\C^r$ preserves these
formulas.  Let $H_n=H_{m,n}^2\otimes\C^r$ and let $P_n$ denote its
orthogonal projection.

\begin{lemma}\label{op:graded}
A closed subspace $\mathcal M\subset H_m^2\otimes\C^r$ invariant under
the coordinates and under $U_t f(z)=f(e^{it}z)$ is the closure of a
finitely generated homogeneous submodule $M\subset R\otimes\C^r$.
Its polynomial part is exactly $M$.  Conversely, the closure of a
homogeneous polynomial submodule has these properties.
\end{lemma}
\begin{proof}
The circle action is unitary and strongly continuous, as follows first
on polynomials and then by density.  The Fourier projection is the
norm-convergent Bochner integral
\[
 P_nf=\frac1{2\pi}\int_0^{2\pi}e^{-int}U_tf\,dt.
\]
For $f\in\mathcal M$, its Riemann sums belong to $\mathcal M$, so
$P_nf\in\mathcal M\cap H_n$.  The partial sums $\sum_{n=0}^NP_nf$
converge to $f$.  Thus $M=\mathcal M\cap(R\otimes\C^r)$ is homogeneous
and dense in $\mathcal M$.  Coordinate invariance makes it an
$R$-submodule.  Since $R$ is Noetherian and $R\otimes\C^r$ is finite
free, $M$ is finitely generated.  Replacing each generator by its finitely
many homogeneous parts gives homogeneous generators; all those parts
belong to $M$ and their sums recover the original generators.

Conversely, boundedness of the coordinates and unitarity of $U_t$ preserve
both invariances upon taking closure.  If a polynomial $f$ belongs to
$\overline M$, then $P_nf$ lies in the closure of $M_n$ for each $n$.
The latter space is finite dimensional and closed.  Hence every
homogeneous part of $f$ belongs to $M$, and their finite sum does too.
\end{proof}

Write $Q=\mathcal M^\perp=\widehat\bigoplus_{n\ge0}Q_n$.
Invariance of $\mathcal M$ implies $Z_j^*Q\subset Q$.  Indeed,
$\langle Z_j^*q,u\rangle=\langle q,Z_ju\rangle=0$ for
$q\in Q$ and $u\in\mathcal M$.  Thus the quotient tuple
$S_j=P_QZ_j|_Q$ has adjoint $Z_j^*|_Q$.

\subsection{Compression identities}
The degree-$n$ defect column is
\[
 B_n:Q_n\longrightarrow M_{n+1}^{\oplus m},\qquad
 B_nq=(B_{1,n}q,\ldots,B_{m,n}q),
\]
so $A_n=B_n^*B_n$.  Its direct sum is a bounded operator because
each coordinate is bounded.

\begin{lemma}\label{op:commutators}
For $n\ge0$,
\begin{align}
 [Z_i^*,Z_j]|_{H_n}
   &=\frac{\delta_{ij}I-Z_jZ_i^*}{n+1}\bigg|_{H_n},
       \label{eq:op-ambient}\\
 [S_i^*,S_j]|_{Q_n}
   &=P_{Q_n}[Z_i^*,Z_j]|_{Q_n}-B_{i,n}^*B_{j,n}.
       \label{eq:op-quotient}
\end{align}
For $n\ge1$, the restriction tuple satisfies
\begin{equation}\label{eq:op-restriction}
 [V_i^*,V_j]|_{M_n}
 =P_{M_n}[Z_i^*,Z_j]|_{M_n}+B_{j,n-1}B_{i,n-1}^*.
\end{equation}
The last term is zero in degree zero.  In particular, a bound
$A_n\le C(n+1)^{-1}I$ gives $O((n+1)^{-1})$ bounds for all
quotient and restriction commutator blocks.
\end{lemma}
\begin{proof}
For $n>0$, apply \eqref{eq:op-shift} to a homogeneous polynomial:
\[
 Z_i^*Z_jf=\frac{\delta_{ij}f+z_j\partial_if}{n+1},\qquad
 Z_jZ_i^*f=\frac{z_j\partial_if}{n}.
\]
Their difference gives \eqref{eq:op-ambient}; constants give the same
identity in degree zero.  Relative to $\mathcal M\oplus Q$, write
\[
 Z_j=\begin{pmatrix}V_j&B_j\\0&S_j\end{pmatrix}.
\]
Multiplying this matrix and the adjoint matrix for $Z_i$ gives
\eqref{eq:op-quotient} from the lower-right corner and
\eqref{eq:op-restriction} from the upper-left corner.  Grading specifies
the displayed incoming and outgoing indices.

Since $B_{j,n}^*B_{j,n}\le A_n$, each
$\|B_{i,n}^*B_{j,n}\|$ is at most $\|A_n\|$.  The ambient block
has norm at most $2/(n+1)$.  The same argument in incoming degree
$n-1$ gives the restriction bound, after increasing a constant to include
degree zero.
\end{proof}

The defect is an operator-norm question.  Its trace alone does not give
the needed bound in higher dimension.  In fact an exact identity is
available without any division estimate.

\begin{proposition}\label{op:trace}
If $h_n=\dim Q_n$, then
\begin{equation}\label{eq:op-trace}
 \Tr A_n=\frac{n+m}{n+1}h_n-h_{n+1}.
\end{equation}
\end{proposition}
\begin{proof}
The forward weights in \eqref{eq:op-shift} give
$\sum_jZ_j^*Z_j|_{H_n}=(n+m)(n+1)^{-1}I$.
Orthogonally splitting each forward product into its ideal and quotient
parts yields
\[
 A_n=\frac{n+m}{n+1}I_{Q_n}-\sum_j S_{j,n}^*S_{j,n}.
\]
On positive degrees the backward weights give
$\sum_jZ_jZ_j^*=I$.  Since $Q$ is coinvariant,
$\sum_jS_{j,n}S_{j,n}^*=I_{Q_{n+1}}$.  Taking finite-dimensional
traces proves the formula.
\end{proof}

Suppose that the quotient is infinite dimensional, so that its support
dimension satisfies $1\leq\delta\leq m$, and write $a>0$ for the leading
coefficient of its eventual Hilbert polynomial (of degree $\delta-1$
by Hilbert--Serre; see Lemma~\ref{op:hilbert}). Rewriting
\eqref{eq:op-trace} gives
\begin{equation}\label{eq:op-trace-asymptotic}
 \Tr A_n=\frac{m-1}{n+1}h_n-(h_{n+1}-h_n)
       =(m-\delta)a n^{\delta-2}+O(n^{\delta-3}).
\end{equation}
Indeed, if $\delta\geq2$, the polynomial difference has leading term
$(\delta-1)a n^{\delta-2}$ and remainder $O(n^{\delta-3})$; if
$\delta=1$, the difference is eventually zero and the same expansion
holds. Thus $\Tr A_n\sim(m-\delta)a n^{\delta-2}$ when $\delta<m$.
When $\delta=m$ the leading coefficient vanishes; the trace is then
$O(n^{m-3})$ and may be identically zero, as for $M=0$. In all cases the
average eigenvalue, $\Tr A_n/h_n$, is $O(n^{-1})$. For $\delta=1$, $h_n=a$
eventually, so positivity gives the full estimate
$\|A_n\|\leq (m-1)a/(n+1)$ for large $n$. This is the dimension-one
trace argument of Guo and Wang for finite coefficient multiplicity
\cite[Proposition~4.1, p.~920]{GuoWang}. In higher support dimension the
trace estimate generally permits concentration in a small subspace.
The lifting argument below supplies the required uniform operator-norm
bound without a restriction on the support dimension.

\subsection{Counting singular values in the original grading}
For a compact operator $T$, let $s_k(T)$ be its singular values in
nonincreasing order, starting at $k=1$.  For $p>0$, write
\[
 T\in\cS_{p,\infty}\quad\Longleftrightarrow\quad
 \sup_{k\ge1} k^{1/p}s_k(T)<\infty.
\]

\begin{lemma}\label{op:counting}
Suppose $T=\bigoplus_{n\ge0}T_n$ has mutually orthogonal domain and
range blocks, with
\[
 \dim\operatorname{dom}T_n\le C_0(n+1)^{\delta-1},\qquad
 \|T_n\|\le C_1(n+1)^{-a},
\]
where $a,\delta>0$.  Then $T$ is compact,
$T\in\cS_{\delta/a,\infty}$, and $T\in\cS_p$ for $p>\delta/a$.
\end{lemma}
\begin{proof}
Finite block truncations are finite rank and converge in norm because
the block norms tend to zero.  If a singular value of $T_n$ exceeds
$t>0$, then $n+1<(C_1/t)^{1/a}$.  The counting function therefore obeys
\[
 \#\{k:s_k(T)>t\}
 \le\sum_{n+1<(C_1/t)^{1/a}}C_0(n+1)^{\delta-1}
 \le C_2t^{-\delta/a}
\]
for small $t$, which is the weak Schatten assertion.  Alternatively,
for $p>\delta/a$,
\[
 \sum_n\|T_n\|_{\cS_p}^p
 \le C_0C_1^p\sum_n(n+1)^{\delta-1-ap}<\infty.
\]
\end{proof}

\begin{lemma}\label{op:hilbert}
Let $T=(R\otimes\C^r)/M$.  If $T$ is infinite dimensional and
$\delta=\dim\Supp T$, then
\[
 \dim Q_n=\dim T_n=O((n+1)^{\delta-1}).
\]
If $\delta=0$, then $T$ is finite dimensional.
\end{lemma}
\begin{proof}
In degree $n$, orthogonal projection gives an isomorphism
$T_n\simeq Q_n$ of finite-dimensional vector spaces.  The
Hilbert--Serre theorem states that the Hilbert function of a nonzero
finite graded module of positive support dimension is eventually a
polynomial of degree $\delta-1$, with positive leading coefficient.
Equivalently, its Hilbert series has a pole of order $\delta$ at $1$.
This applies to $T$ because it is a quotient of a finite free standard
graded module.  The finitely many earlier degrees are absorbed into a
constant.  In dimension zero its Hilbert function is eventually zero,
so its total vector-space dimension is finite; see
\cite{SerreFAC,StacksHilbert}.
\end{proof}

Applying Lemma~\ref{op:counting} with $a=1$ to
\eqref{eq:main-commutator} proves the quotient assertions of
Theorem~\ref{thm:main}, once \eqref{eq:main-degree} is established.
For the column $B$, use $\|B_n\|^2=\|A_n\|$ and $a=1/2$.
The commutator $[P_{\mathcal M},Z_j]$ is the block $B_j$ extended by
zero, so it has the same nonzero singular values.  For the restriction
tuple use \eqref{eq:op-restriction} and
\[
 \dim M_n\le r\binom{n+m-1}{m-1}=O((n+1)^{m-1}).
\]
This proves its weak exponent $m$.  If $T$ is finite dimensional then
$Q$ is finite dimensional, and the assertions involving $Q$ have finite
rank.  Notice that the restriction exponent and the quotient exponent
come from different dimension counts.

% END sections/operators.tex


% BEGIN sections/projective.tex
\section{Step 1: Projective norms and the lifting problem}
\label{sec:projective}

We now treat a scalar homogeneous ideal.  This section keeps the relation
between the holomorphic Hilbert space and the geometric $L^2$ spaces
explicit.  In particular, the $\bar\partial$ equations used later are
equations for bundle-valued forms on projective space.

\subsection{Saturation changes only finitely many blocks}
Let $I\subset R=\C[z_1,\ldots,z_m]$ be homogeneous and let
$\mathfrak m=(z_1,\ldots,z_m)$.  Put
\[
 I^{\mathrm{sat}}=\bigcup_{k\ge0}(I:\mathfrak m^k).
\]

\begin{lemma}\label{geo:saturation}
The graded module $I^{\mathrm{sat}}/I$ is finite dimensional.
Consequently $I_n=I_n^{\mathrm{sat}}$ for all sufficiently large $n$,
and a bound $A_n\le C(n+1)^{-1}I$ for the saturation gives the same
type of bound for $I$ after increasing $C$.
\end{lemma}
\begin{proof}
The ascending chain of colon ideals stabilizes, since $R$ is Noetherian.
Thus $\mathfrak m^k I^{\mathrm{sat}}\subset I$ for one $k$.
The quotient $I^{\mathrm{sat}}/I$ is finitely generated and is annihilated
by $\mathfrak m^k$, so it is a finite module over the finite-dimensional
algebra $R/\mathfrak m^k$.  It has only finitely many nonzero degree
spaces.  In high degree the orthogonal projections onto $I_n$ and onto
its saturation are literally equal in the original inner product.
The defect in degree $n$ uses only degrees $n$ and $n+1$, so those
defects coincide eventually.  There are finitely many remaining
finite-dimensional blocks and their norms can be absorbed into $C$.
\end{proof}

For example, in $\C[x,y]$ the ideal $I=(x^2,xy)$ has saturation $(x)$.
Indeed $\mathfrak m x\subset I$, and $(x)$ is saturated since
$\C[x,y]/(x)\simeq\C[y]$ has no $\mathfrak m$-torsion. The quotient
$(x)/I$ is the one-dimensional space spanned by the class of $x$.
In this example the projective scheme is the reduced point $[0:1]$;
the finite-length discrepancy occurs only at the cone vertex. In general,
saturation preserves the entire projective ideal sheaf, including any
nonreduced structure supported on projective space.

For $m=1$, a scalar homogeneous ideal is $0$ or $(z_1^a)$ for some
$a\ge0$.  The first has zero defect; the others have finite-dimensional
quotient.  Likewise, $I^{\mathrm{sat}}=R$ gives a finite quotient in
any dimension.  We may therefore assume $m\ge2$ and that the saturated
ideal sheaf is nonzero and proper.  The zero ideal has zero defect.

Write $N=m-1$, $W=\Pj^N$, and $L=\cO_W(1)$.  The saturation determines
an algebraic coherent ideal sheaf $\cI\subset\cO_W$, whose
analytification is also denoted $\cI$.  In nonnegative degree,
\begin{equation}\label{eq:geo-sections}
 I^{\mathrm{sat}}_n=H^0(W,\cI\otimes L^n)\subset H^0(W,L^n)=R_n.
\end{equation}
For completeness, an algebraic section represented by $p\in R_n$
belongs to $\cI(n)$ precisely when each standard localization places
it in the localized ideal.  Thus for each $i$ some $z_i^{k_i}p$ lies
in $I$, with $k_i$ chosen positive.  Every monomial of degree greater than
$\sum_i(k_i-1)$ is divisible by one of these powers, so
$\mathfrak m^kp\subset I$ for a sufficiently large $k$.
This is saturation.  The converse is immediate by localization.
The equality with holomorphic global sections follows from GAGA.
Analytification also preserves exact sequences of coherent sheaves;
these are the exactness and global-section comparisons used below
\cite{SerreGAGA}.

\subsection{The original norm as a projective integral}
Give $L$ its Fubini--Study metric and normalize the Fubini--Study volume
$dV$ on $W$ to have total mass one.  A degree-$n$ homogeneous polynomial
$p$ defines a section of $L^n$ with pointwise norm
\[
 |p([z])|_{L^n}^2=\frac{|p(z)|^2}{|z|^{2n}}.
\]
The quotient is unchanged under $z\mapsto\lambda z$.

\begin{lemma}\label{geo:norm}
For $p\in R_n$,
\begin{equation}\label{eq:geo-norm}
 \|p\|_{\DA,n}^2=c_n\|p\|_{\FS,n}^2,
 \qquad
 c_n=\frac{(n+m-1)!}{(m-1)!\,n!}.
\end{equation}
Moreover,
\begin{equation}\label{eq:geo-ratio}
 \frac{c_n}{c_{n+1}}=\frac{n+1}{n+m}.
\end{equation}
Orthogonal projections within a fixed degree are the same for the two
inner products.
\end{lemma}
\begin{proof}
The Hopf projection pushes normalized sphere measure $d\sigma$ on
$S^{2m-1}$ to $dV$.  Hence the projective squared norm of $p$ is
$\int_{S^{2m-1}}|p|^2d\sigma$.  For $|\alpha|=n$, the Gaussian
integral factors as
\[
 \int_{\C^m}|z^\alpha|^2e^{-|z|^2}\,d\lambda(z)=\pi^m\alpha!.
\]
In polar coordinates its radial factor, including the area of the unit
sphere, is $\pi^m(n+m-1)!/(m-1)!$.  Therefore
\[
 \|z^\alpha\|_{\FS,n}^2
 =\frac{(m-1)!\,\alpha!}{(n+m-1)!}.
\]
Coordinatewise rotations make distinct monomials orthogonal for this
norm.  Comparison with $\alpha!/n!$ proves
\eqref{eq:geo-norm}; taking the ratio gives
\eqref{eq:geo-ratio}.  A scalar change of inner product leaves
orthogonal projections unchanged.
\end{proof}

\begin{remark}\label{geo:dbar-space}
The Drury--Arveson norm is not an ordinary positive-measure norm on the
closed ball when $m\ge2$.  Such a measure would have mass
$\|1\|^2=1$, but would also satisfy
$\int|z|^2d\mu=\sum_j\|z_j\|^2=m>1$.
Lemma~\ref{geo:norm} instead identifies its homogeneous degree spaces
with $H^0(W,L^n)$; for $f=\sum_{n\geq0}f_n$ its full squared norm is
\[
 \|f\|_{\DA}^2=\sum_{n\geq0}c_n\|f_n\|_{\FS,n}^2.
\]
The geometric ambient spaces used in this proof are
$L^2_{0,j}(W,E_q\otimes L^n)$ for finite-rank coefficient bundles
$E_q$.  Intermediate solutions need not be holomorphic.  After the
correction, a holomorphic section of a nonnegative power of $L$ is an
actual homogeneous polynomial, to which \eqref{eq:geo-norm} applies.
No $\bar\partial$ operator on a space consisting only of holomorphic
functions is being used.
\end{remark}

\subsection{The geometric task}
The scalar estimate will be recovered in Section~\ref{sec:completion}
from the compatible metrics of Theorem~\ref{thm:D} and the lifting bound
of Theorem~\ref{geo:lifting}. In particular, if $\Phi a=F$ is a
holomorphic lift of $F\in I_{n+1}$ through a fixed complete row of ideal
generators, then
\[
 P_{Q_n}\partial_{z_j}F=P_{Q_n}((\partial_{z_j}\Phi)a)
\]
for $n$ beyond the saturation discrepancy. This follows from the ordinary
product rule. To control the right side it is enough to construct a
lift and a differentiated-row estimate in one compatible coefficient
metric. The rest of Step 1 constructs the algebraic presentation on which
those estimates will be based.

% END sections/projective.tex


% BEGIN sections/algebra.tex
\section{Complete presentations and their normalized limits}
\label{sec:algebra}

Throughout this section, \(W=\mathbb P^N\), \(L=\mathcal O_W(1)\), and
\(R=\mathbb C[Z_0,\ldots,Z_N]\). We assume \(N\geq1\) and fix a nonzero
proper coherent algebraic ideal sheaf \(\mathscr I\subset\mathcal O_W\).
The zero and unit ideals, and the one-variable operator problem, require no
presentation of the kind constructed below.

The local estimates will involve changes of center, changes of scale, and
possibly ill-conditioned constant changes of basis. An arbitrary
finite set of generators is poorly suited to taking these limits. We
instead use complete spaces of sections in degrees chosen uniformly over
proper parameter spaces. This choice controls the dimensions of the
polynomial images in every degree, including at a limiting matrix.

\paragraph{Resolution and the choice of degrees.}
The existence of a partial locally free resolution by sufficiently
positive evaluation spaces follows from the classical generation
and regularity theorems. In the context of essential normality,
\cite{DouglasJabbariTangYu} gives a resolution by analytic Hilbert
modules for monomial ideals; in polynomial membership problems,
\cite{AnderssonWulcan} uses graded free resolutions and projective
Dolbeault complexes. Our choice here serves a different quantitative
purpose. A presentation must still describe the correct image sheaves
when affine dilations and coefficient normalizations degenerate. The
entire proper quotient family, not just the initial fiber, determines
the generating degree. The uniformity proved in this section is the
input to the fixed-scale division theorem.

\subsection{Regularity and complete evaluation maps}

For a coherent sheaf \(\mathscr F\) on \(W\), write
\(\mathscr F(t)=\mathscr F\otimes L^t\), and let
\[
 P_{\mathscr F}(t)=\chi(W,\mathscr F(t))
\]
be its Hilbert polynomial. We use the convention
\(R(-D)_\ell=R_{\ell-D}\), with \(R_j=0\) for \(j<0\).

Recall that \(\mathscr F\) is \(a\)-regular if
\[
 H^i(W,\mathscr F(a-i))=0,\qquad 1\leq i\leq N.
\]
We will use the following precise consequences of
Castelnuovo--Mumford regularity.

\begin{lemma}[Complete sections in a regular degree]
\label{alg:regular-sections}
If \(\mathscr F\) is \(a\)-regular, then for every \(t\geq a\):
\begin{enumerate}
\item \(\mathscr F\) is \(t\)-regular, and \(\mathscr F(t)\) is generated
by its global sections;
\item \(H^i(W,\mathscr F(t))=0\) for \(i>0\), so that
\(h^0(W,\mathscr F(t))=P_{\mathscr F}(t)\);
\item multiplication induces a surjection
\[
 H^0(W,L)\otimes H^0(W,\mathscr F(t))
 \longrightarrow H^0(W,\mathscr F(t+1)).
\]
\end{enumerate}
Consequently, for \(b\geq a\) and \(\ell\geq b\), the map
\[
 R_{\ell-b}\otimes H^0(W,\mathscr F(b))
 \longrightarrow H^0(W,\mathscr F(\ell))
\]
is surjective.
\end{lemma}

\begin{proof}
Persistence of regularity, generation, and the multiplication assertion
are the Castelnuovo--Mumford regularity lemmas
\cite[Tags 08A6--08A8]{StacksRegularity}. Persistence up to \(t+i\),
followed by the defining vanishing in cohomological degree \(i\), gives
\(H^i(W,\mathscr F(t))=0\). The equality with the Hilbert polynomial is
then the Euler characteristic formula. Iterating multiplication gives a
surjection from the corresponding tensor power of \(H^0(W,L)\).
Multiplication of linear forms is symmetric, so this map factors through
\(R_{\ell-b}\), proving the last assertion.
\end{proof}

\begin{lemma}[One degree for a whole quotient family]
\label{alg:uniform-regularity}
Fix \(N,r,D\) with \(r>0\), and a polynomial \(P\). There is an integer
\(a\), depending only on these data, such that the kernel of every
quotient
\[
 \mathcal O_W(-D)^r\twoheadrightarrow\mathscr F,
 \qquad P_{\mathscr F}=P,
\]
is \(a\)-regular.
\end{lemma}

\begin{proof}
Twisting by \(D\) gives a quotient of \(\mathcal O_W^r\) whose Hilbert
polynomial is \(P(t+D)\). The uniform regularity theorem for kernels of
such quotients gives a regularity bound depending only on \(N,r\), and
\(P(t+D)\) \cite[Tag 08AG]{StacksRegularity}. Untwisting adds \(D\) to
that bound.
\end{proof}

The degree in this lemma is chosen for the entire quotient family,
including fibers with additional torsion or embedded structure. It is
not chosen only for the particular sheaf with which the argument begins.

\begin{proposition}[Complete presentation through dimension \(N\)]
\label{alg:complete}
There are integers
\[
 0\leq D_0\leq D_1\leq\cdots\leq D_N
\]
and finite-dimensional spaces \(V_q\) giving a complex
\begin{equation}
 E_N\xrightarrow{d_N}\cdots\xrightarrow{d_2}E_1
 \xrightarrow{d_1}E_0\xrightarrow{d_0=\Phi}\mathcal O_W,
 \qquad E_q=V_q\otimes L^{-D_q},
 \label{alg:complex}
\end{equation}
with the following properties. Put \(E_{-1}=\mathcal O_W\),
\(D_{-1}=0\), \(r_q=\dim V_q\), \(r_{-1}=1\), and
\[
 \mathscr K_0=\mathscr I,\qquad
 \mathscr K_q=\ker d_{q-1}\quad(q\geq1).
\]
Then
\[
 \operatorname{im}d_q=\mathscr K_q,
 \qquad
 V_q=H^0(W,\mathscr K_q(D_q)).
 \tag{\(\ast\)}
\]
Moreover, \(D_0\) is a regular degree for every ideal sheaf with quotient
Hilbert polynomial \(P_{\mathcal O_W/\mathscr I}\); for \(q\geq1\),
\(D_q\) is a regular degree for every kernel in
\[
 \operatorname{Quot}_{W}(E_{q-1},P_{\mathscr K_{q-1}}).
 \tag{\(\ast\ast\)}
\]
Every nonzero matrix \(d_q\) has columns linearly independent over
constant scalars. Its entries are homogeneous polynomials of degree
\(e_q=D_q-D_{q-1}\).
\end{proposition}

\begin{proof}
Set \(P_0=P_{\mathscr I}\) and
\(P_X=P_{\mathcal O_W/\mathscr I}\). Apply
Lemma~\ref{alg:uniform-regularity} to the quotients of
\(\mathcal O_W\) with polynomial \(P_X\), and choose a common regular
degree \(D_0\geq0\). The complete evaluation map
\[
 H^0(W,\mathscr I(D_0))\otimes L^{-D_0}
 \longrightarrow \mathscr I
\]
is surjective by Lemma~\ref{alg:regular-sections}. Its composition with
\(\mathscr I\hookrightarrow\mathcal O_W\) is \(d_0=\Phi\).

Suppose \(E_{q-1}\) and \(d_{q-1}\) have been constructed. Their kernel
fits into an exact sequence
\begin{equation}
 0\longrightarrow\mathscr K_q\longrightarrow E_{q-1}
 \longrightarrow\mathscr K_{q-1}\longrightarrow0.
 \label{alg:kernel-sequence}
\end{equation}
Write \(P_q=P_{\mathscr K_q}\). Additivity gives
\begin{equation}
 P_q(t)=r_{q-1}\binom{t-D_{q-1}+N}{N}-P_{q-1}(t).
 \label{alg:polynomial-recursion}
\end{equation}
The binomial is interpreted as a polynomial in \(t\). The source
\(E_{q-1}\) and the quotient polynomial \(P_{q-1}\) are now fixed.
Choose \(D_q\geq D_{q-1}\) regular for every kernel in
\((\ast\ast)\), and put
\[
 V_q=H^0(W,\mathscr K_q(D_q)),\qquad
 E_q=V_q\otimes L^{-D_q}.
\]
The complete evaluation onto \(\mathscr K_q\), followed by its inclusion
in \(E_{q-1}\), defines \(d_q\). Its image is precisely the preceding
kernel. If a kernel is zero, all subsequent bundles may be taken to be
zero.

Choose a basis of each nonzero \(V_q\). The columns of \(d_q\) are exactly
these basis sections inside
\(H^0(W,E_{q-1}(D_q))\), so a constant linear relation between the
columns is zero only when its coefficients vanish. Since all summands
of the source and target have common twists, the entries have degree
\(D_q-D_{q-1}\). This completes the induction.
\end{proof}

The presentation need not terminate with an injective map \(d_N\). The
analytic correction will use forms only through antiholomorphic degree
\(N\), so exactness at \(E_0,\ldots,E_{N-1}\) is sufficient. Nor do we
assert that the associated graded polynomial complex is exact in small
degrees.

\paragraph{The last map need not be injective.}
For example, on $\Pj^1$ take $\mathscr I=(Z_0)\simeq\cO(-1)$,
$D_0=2$ and $D_1=4$. These degrees meet the uniform regularity
requirements: the ideals of length-one subschemes are isomorphic to
$\cO(-1)$, and the kernels in the next Quot family are torsion-free
rank-one sheaves of degree $-3$, hence isomorphic to $\cO(-3)$.
The complete evaluations are
\[
 d_0=(Z_0^2, Z_0Z_1),\qquad
 d_1=\begin{pmatrix}-Z_1\\ Z_0\end{pmatrix}(Z_0,Z_1).
\]
Here $E_0=\cO(-2)^2$ and $E_1=\cO(-4)^2$. The row $d_0$ has
kernel $\cO(-3)$: the map $(-Z_1,Z_0)^{\mathsf T}$ identifies
that line bundle with the kernel, because the two coordinates
have no common projective zero. The row $(Z_0,Z_1)$ is a surjection
$\cO(-4)^2\to\cO(-3)$ obtained from the complete section space
of $\cO(1)$. Its kernel is $\cO(-5)$ by the same calculation.
Thus $\im d_1=\ker d_0$ but $\ker d_1\simeq\cO(-5)\ne0$.
This is a partial presentation of the permitted kind, not a complex
asserted exact at its left endpoint.

\begin{lemma}[The actual graded image]
\label{alg:evaluation-image}
Let \(J_q\) be the image of the homogeneous polynomial matrix
\[
 R(-D_q)^{r_q}\longrightarrow R(-D_{q-1})^{r_{q-1}}
\]
representing \(d_q\). Then
\begin{equation}
 (J_q)_\ell=
 \begin{cases}
  0,&\ell<D_q,\\
  H^0(W,\mathscr K_q(\ell)),&\ell\geq D_q,
 \end{cases}
 \qquad
 \dim(J_q)_\ell=
 \begin{cases}
  0,&\ell<D_q,\\
  P_q(\ell),&\ell\geq D_q.
 \end{cases}
 \label{alg:image-function}
\end{equation}
The same formulas hold for the complete evaluation of every kernel in
the corresponding parameter family.
\end{lemma}

\begin{proof}
Below \(D_q\), the source homogeneous component vanishes. At degree
\(D_q\), its image is the entire section space by construction. In every
higher degree, Lemma~\ref{alg:regular-sections} identifies the polynomial
multiples of these columns with all sections of
\(\mathscr K_q(\ell)\). Including those sections in
\(H^0(W,E_{q-1}(\ell))\) gives the stated equality. The dimension is its
Hilbert polynomial because all higher cohomology vanishes. The same
regularity bound was imposed on every member of the family.
\end{proof}

\subsection{Sampled Gram forms and compact normalization}

Use standard projective affine charts and their standard line frames.
For a chart \(\alpha\), let \(d_{q,\alpha}(w)\) be the polynomial
coefficient matrix of \(d_q\). Centers \(p\) range in fixed relatively
compact subsets of these charts. Choose \(\varepsilon_0>0\) so that all
balls used in the following definition remain in the corresponding
chart whenever \(0<\varepsilon\leq\varepsilon_0\). Let \(\mu\) be
probability Lebesgue measure on the unit ball \(B_1\subset\mathbb C^N\).

For \(v\in V_q\), define
\begin{equation}
\begin{split}
 |v|_{G_{q,\alpha,\varepsilon}(p)}^2
 =\int_{(B_1)^{q+1}}
 \bigl|d_{0,\alpha}(p+\varepsilon z_0)\cdots
 d_{q,\alpha}(p+\varepsilon z_q)v\bigr|^2
 \prod_{j=0}^{q}d\mu(z_j).
\end{split}
\label{alg:gram-product}
\end{equation}
The sample points in the different factors are independent. Sampling
all factors at the same point would give zero whenever two consecutive
maps occur, because their composition is zero.

\begin{lemma}[Positivity and recursion]
\label{alg:gram}
The forms \(G_{q,\alpha,\varepsilon}(p)\) are smooth and positive
definite at every nonzero stage. With \(G_{-1,\alpha,\varepsilon}=1\),
they satisfy the frozen-center recursion
\begin{equation}
 G_{q,\alpha,\varepsilon}(p)=
 \int_{B_1}d_{q,\alpha}(p+\varepsilon z)^*
 G_{q-1,\alpha,\varepsilon}(p)
 d_{q,\alpha}(p+\varepsilon z)\,d\mu(z).
 \label{alg:gram-recursion}
\end{equation}
\end{lemma}

\begin{proof}
Integrate first over \(z_0,\ldots,z_{q-1}\) in
\eqref{alg:gram-product}; the result is exactly
\eqref{alg:gram-recursion}. Inductively, if its value on \(v\) is zero,
the nonnegative continuous integrand vanishes on \(B_1\). Positivity of
the preceding form implies
\[
 d_{q,\alpha}(p+\varepsilon z)v=0\qquad(z\in B_1).
\]
Since \(\varepsilon>0\), this is a polynomial identity on a nonempty
open set. It is a constant linear relation between the columns of
\(d_q\), which are independent by Proposition~\ref{alg:complete}.
Therefore \(v=0\). Smoothness follows directly from integration of a
polynomial in \(p,\overline p\), and \(\varepsilon\) over a fixed compact
domain.
\end{proof}

For a fixed center and scale, abbreviate
\[
 R_q=G_{q,\alpha,\varepsilon}(p)^{1/2},
 \qquad R_{-1}=1,
\]
and define
\begin{equation}
 A_q(\zeta)
 =
 R_{q-1}\,d_{q,\alpha}(p+\varepsilon\zeta)\,R_q^{-1}.
 \label{alg:normalized}
\end{equation}
The \(R_q\) in this formula are constant matrices: the center is frozen
while \(\zeta\) varies.

\begin{thmB}[Complete presentations stable under normalized limits]\label{thm:B}
The complete presentation of Proposition~\ref{alg:complete}, equipped
with the Gram forms \eqref{alg:gram-product}, has the following properties.
The tuples \eqref{alg:normalized}, as the center ranges over the relatively
compact sets of a finite buffered affine atlas and $0<\varepsilon\le
\varepsilon_0$, have compact coefficient closure. Homogenize each map to
its prescribed degree $e_q=D_q-D_{q-1}$. In the coordinates of the corresponding chart (or after the fixed
coordinate permutation identifying that chart with $\{Z_0\neq0\}$), each
simultaneous limit is a complex, exact at $E_0,\ldots,E_{N-1}$, with
image of its first row an ideal sheaf of Hilbert polynomial $P_0$.
Since the atlas is finite, a convergent sequence may be analyzed on a
subsequence in a single chart. For every normalized or limiting
map, its homogeneous polynomial image $J'_q$ satisfies
\[
 \dim (J'_q)_\ell=\begin{cases}0,&\ell<D_q,\\
 P_q(\ell),&\ell\ge D_q.\end{cases}
\]
Here $P_q$ is the polynomial of $\mathscr K_q$ and the degrees in the
image are measured with the prescribed source and target shifts.
\end{thmB}
The following results prove, separately, the coefficient compactness,
retention of complete sections, exactness, and all-degree image formula.

\begin{proposition}[Compactness of normalized tuples]
\label{alg:compact}
The tuples \((A_0,\ldots,A_N)\) have compact closure in the
finite-dimensional spaces of their polynomial coefficients. At every
nonzero stage,
\begin{equation}
 \int_{B_1}A_q(\zeta)^*A_q(\zeta)\,d\mu(\zeta)=I_{r_q}.
 \label{alg:orthonormal-columns}
\end{equation}
Every simultaneous limiting tuple has linearly independent columns at
each such stage and satisfies
\[
 A_{q-1}^{\infty}(\zeta)A_q^{\infty}(\zeta)=0.
\]
\end{proposition}

\begin{proof}
Substitution into the integral in
\eqref{alg:orthonormal-columns} gives
\[
 R_q^{-1}
 \left(\int_{B_1}d_{q,\alpha}(p+\varepsilon\zeta)^*
 G_{q-1,\alpha,\varepsilon}(p)
 d_{q,\alpha}(p+\varepsilon\zeta)\,d\mu(\zeta)\right)
 R_q^{-1}
 =I.
\]
The columns thus form an orthonormal frame in the fixed
finite-dimensional space of \(\mathbb C^{r_{q-1}}\)-valued polynomials
of degree at most \(e_q\), with its probability-ball inner product.
The space of such frames is compact. Take its product over the finitely
many stages. Orthonormality survives coefficient limits, so no column
dependence appears. The intermediate factors \(R_{q-1}^{-1}R_{q-1}\)
cancel in a product of consecutive normalized matrices. The original
composition is zero, and the resulting polynomial identity survives
every simultaneous limit.
\end{proof}

Compactness and independence of the columns alone would not control
their polynomial multiples. For example, the homogeneous row
\[
 (X^2,\ XY+tY^2)
\]
has two independent columns even at \(t=0\), but its degree-three image
has dimension four for \(t\neq0\) and dimension three at zero. The
complete presentations will rule out this loss of image dimension.

\begin{example}[Why regularity must be uniform over the whole family]
\label{alg:degeneration-test}
On \(\mathbb P^2\), with coordinates \([Z:X:Y]\), consider
\[
 I_s=(X^2,XY,XZ-sY^2),\qquad s\ne0.
\]
On \(Z=1\) its quotient is
\(\mathbb C[x,y]/(x-sy^2,y^3)\), and it has no points on \(Z=0\).
Thus it defines a length-three scheme. Its quadratic generators are
its complete quadratic section space. Indeed, restriction of the six
quadratic monomials to the local quotient has image spanned by
\(1,y,y^2\), so the kernel has dimension three.
Nevertheless, the column-independent limit of this quadratic row is
\((X^2,XY,XZ)\), which generates the line ideal sheaf \((X)\).
Its cubic image has dimension six, whereas the preceding cubic images
have dimension seven.

The omitted direction is explicit:
\[
 Y(XZ-sY^2)-Z(XY)=-sY^3.
\]
Adding \(Y^3\) to the family of generators does not change its fibers
for $s\ne0$. The resulting quotient sheaf is supported in the chart
$Z\ne0$, since $X^2$ and $Y^3$ vanish on its support. Its restriction
there corresponds to
$\mathbb C[s,x,y]/(x-sy^2,y^3)\simeq\mathbb C[s,y]/(y^3)$,
which is free over $\mathbb C[s]$. It is therefore flat over the
parameter line, and its ideal kernel is flat as well, by the exact
sequence with the flat ambient structure sheaf. At $s=0$ the ideal
sheaf is $(X,Y^3)$. Its complete degree-two section space generates
only $(X)$, so the ideal sheaf is not generated in degree two and
cannot be two-regular. A common regularity degree for the \emph{ideal
sheaves} of all length-three subschemes must therefore be at least
three. The quotient sheaves themselves have no higher cohomology in
any twist, since their support is zero dimensional. At degree three the complete section
space of the limiting ideal sheaf has dimension seven and retains the \(Y^3\) direction at the
limit. This example arises from the local dilation
\((x,y)\mapsto(sx,sy)\) applied to \(I_1\); it tests precisely the
kind of rescaling used in \eqref{alg:normalized}. Completeness for
one fiber alone is insufficient. The degree choice in
Proposition~\ref{alg:complete} is what also controls its limits.
\end{example}

\subsection{Complete section spaces form a closed family}

Work first in the chart \(Z_0\neq0\). The affine change of variables
\(\zeta\mapsto p+\varepsilon\zeta\) is induced by the projective
automorphism
\begin{equation}
 g_{p,\varepsilon}[Z_0:Z']
 =
 [Z_0:pZ_0+\varepsilon Z'].
 \label{alg:projective-change}
\end{equation}
Its matrix is invertible for \(\varepsilon>0\). Homogenizing an entry
of \(d_{q,\alpha}(p+\varepsilon\zeta)\) to the prescribed degree
\(e_q\) gives its pullback by \(g_{p,\varepsilon}\), even if its
observed affine degree is smaller. Coordinate permutations give the
same description in the other standard charts.

Pullback identifies \(g_{p,\varepsilon}^*L\) with \(L\) through this
homogeneous substitution. The target matrix \(R_{q-1}\) acts by a
global automorphism of \(E_{q-1}\): all its summands have the same
line twist. The source factor \(R_q^{-1}\) changes the basis without
changing its span. Thus the homogenized columns of \(A_q\) span
\begin{equation}
 H^0(W,\mathscr K_q'(D_q))
 \subset H^0(W,E_{q-1}(D_q)),
 \qquad
 \mathscr K_q'=R_{q-1}g_{p,\varepsilon}^*\mathscr K_q.
 \label{alg:transformed-sections}
\end{equation}
For \(q\geq1\), the quotient \(E_{q-1}/\mathscr K_q'\) has polynomial
\(P_{q-1}\). At \(q=0\), the quotient polynomial is \(P_X\).

\begin{proposition}[Complete sections at every normalized limit]
\label{alg:complete-limit}
Let \((A_0^\infty,\ldots,A_N^\infty)\) be a simultaneous coefficient
limit, and homogenize \(A_q^\infty\) to degree \(e_q\). Its columns
span the complete space
\[
 H^0(W,\mathscr K_q'(D_q))
 \subset H^0(W,E_{q-1}(D_q))
\]
for a kernel \(\mathscr K_q'\) in the parameter family used to choose
\(D_q\). Their evaluation, after untwisting, generates
\(\mathscr K_q'\).
\end{proposition}

\begin{proof}
For \(q=0\), let \(T\) be the Hilbert scheme of quotients of
\(\mathcal O_W\) with polynomial \(P_X\); for \(q\geq1\), let
\[
 T=\operatorname{Quot}_W(E_{q-1},P_{q-1}).
\]
These are proper parameter spaces
\cite{GrothendieckHilbert,StacksQuot}. Put \(E=E_{q-1}\), and write
their universal quotient sequence on \(W\times T\) as
\begin{equation}
 0\longrightarrow\mathscr U
 \longrightarrow\operatorname{pr}_W^*E
 \longrightarrow\mathscr F
 \longrightarrow0.
 \label{alg:universal-kernel}
\end{equation}
The middle term and \(\mathscr F\) are flat over \(T\). The Tor
sequence shows both that \(\mathscr U\) is flat over \(T\) and that
its fiber is the actual kernel of the corresponding quotient.

Let \(\pi:W\times T\to T\). The choice of \(D_q\) gives
\[
 H^i(W,\mathscr U_t(D_q))=0\quad(i>0),\qquad
 h^0(W,\mathscr U_t(D_q))=r_q
\]
for every fiber. The cohomology-and-base-change theorem therefore makes
\[
 \mathscr V=\pi_*(\mathscr U(D_q))
\]
a locally free sheaf of rank \(r_q\), with fiber exactly
\(H^0(W,\mathscr U_t(D_q))\). This is valid over a nonreduced base:
the vanishing of fiber \(H^1\), flatness, and properness are the needed
hypotheses \cite[Theorem~1.1 and Corollary~1.2]{ConradBaseChange};
see also \cite[III, Theorem~12.11]{Hartshorne}.

The inclusion in \eqref{alg:universal-kernel} gives a map
\[
 \mathscr V\longrightarrow
 H^0(W,E(D_q))\otimes\mathcal O_T
\]
whose fiber maps are injective. It is a subbundle injection. Indeed,
at any point choose a nonzero \(r_q\) by \(r_q\) minor of the fiber
matrix. Its determinant is invertible on a neighborhood, and projecting
onto those coordinates and inverting the minor gives a local left
inverse. This argument also applies to nonreduced bases.

The subbundle defines a morphism
\[
 T\longrightarrow
 \operatorname{Gr}\bigl(r_q,H^0(W,E(D_q))\bigr).
\]
Its image is closed by properness. It is therefore closed in the usual
complex topology as well. Equation \eqref{alg:transformed-sections}
places every positive-radius column span in this image.
Proposition~\ref{alg:compact} says that the limiting columns remain
independent, so their span is the limit in this Grassmannian of the
preceding spans. It is consequently the complete section space of
one of the family kernels. That kernel is \(D_q\)-regular, so these
sections generate it. The case \(r_q=0\) is the corresponding zero
subbundle and needs no Grassmannian coordinate calculation.
\end{proof}

The Gram recursion \eqref{alg:gram-recursion} and boundedness of the
original polynomial matrices on the buffered sampling regions give
uniform upper bounds on $G_{q-1}$, hence on $R_{q-1}=G_{q-1}^{1/2}$.
However, $R_{q-1}$ can approach a singular matrix, and its inverse
need not be uniformly bounded. The proof does not require an invertible
limit of the changes of basis: each positive-radius normalization defines
a point of the same proper parameter space, and its column span is
retained by the closed Grassmannian image.

\begin{theorem}[Exactness of every limiting complex]
\label{alg:exact-limit}
The fixed-degree homogenizations of every simultaneous normalized limit
satisfy
\[
 \operatorname{im}A_0^\infty=\mathscr I^\infty
\]
for an ideal sheaf with polynomial \(P_0\) and quotient polynomial
\(P_X\), and
\[
 \operatorname{im}A_q^\infty=\ker A_{q-1}^\infty
 \qquad(1\leq q\leq N).
\]
\end{theorem}

\begin{proof}
The first assertion follows from
Proposition~\ref{alg:complete-limit}: the columns are the complete
regular-degree sections of an ideal in the Hilbert family and generate
that ideal. Suppose inductively that the actual image of
\(A_{q-1}^\infty\) has Hilbert polynomial \(P_{q-1}\).
The exact sequence onto that image gives
\[
 P_{\ker A_{q-1}^\infty}(t)
 =
 r_{q-1}\binom{t-D_{q-1}+N}{N}-P_{q-1}(t)
 =
 P_q(t).
 \]
Independently, Proposition~\ref{alg:complete-limit} identifies
\(\operatorname{im}A_q^\infty\) with a kernel from the corresponding
Quot family, also of polynomial \(P_q\).

The composition identity in Proposition~\ref{alg:compact},
homogenized to the sum of the two prescribed degrees, gives an
inclusion
\[
 \operatorname{im}A_q^\infty\subseteq\ker A_{q-1}^\infty.
\]
Their coherent quotient has zero Hilbert polynomial. Such a quotient
must vanish. To see this directly, a sufficiently positive twist of
any nonzero coherent sheaf is generated by global sections and has
zero higher cohomology. A nonzero stalk then forces a nonzero global
section, whereas zero Hilbert polynomial would force its space of
sections to have dimension zero. This is a contradiction.
Thus the inclusion is equality, completing the induction.
\end{proof}

Analytification preserves these exact sequences of coherent algebraic
sheaves \cite[Proposition~10]{SerreGAGA}. Accordingly, for
$1\leq q\leq N$, a holomorphic section of $E_{q-1}$ in the kernel
of $d_{q-1}$ (or of the corresponding limiting map) has a holomorphic
preimage under $d_q$ (respectively the limiting $q$th map) locally at
every point of each affine chart. The first row generates its
corresponding analytic ideal sheaf. These assertions supply membership
in the relevant analytic image sheaves;
the estimates for such preimages are established in the next section.

\begin{theorem}[The Hilbert functions do not change at normalized limits]
\label{alg:hilbert}
For every normalized matrix, including every matrix in a simultaneous
coefficient limit, let \(J_q'\) denote its homogeneous polynomial image
from \(R(-D_q)^{r_q}\) to \(R(-D_{q-1})^{r_{q-1}}\). Then, for every
integer \(\ell\),
\begin{equation}
 \dim(J_q')_\ell=
 \begin{cases}
  0,&\ell<D_q,\\
  P_q(\ell),&\ell\geq D_q.
 \end{cases}
 \label{alg:limit-hilbert}
\end{equation}
\end{theorem}

\begin{proof}
Proposition~\ref{alg:complete-limit} gives the complete degree-\(D_q\)
evaluation of a uniformly regular family kernel at each limit.
Lemma~\ref{alg:evaluation-image} applies to that evaluation and
gives \eqref{alg:limit-hilbert}. The same lemma applies to every
positive-radius member.
\end{proof}

\begin{proof}[Proof of Theorem~\ref{thm:B}]
Proposition~\ref{alg:complete} constructs the presentation with the
required common generating degrees. Lemma~\ref{alg:gram} proves that
its Gram forms are positive. Proposition~\ref{alg:compact} proves
compactness and the composition identities at a simultaneous limit.
Proposition~\ref{alg:complete-limit} identifies the column span of each
limiting map with the complete section space of a kernel in its proper
parameter family. Theorem~\ref{alg:exact-limit} proves exactness, and
Theorem~\ref{alg:hilbert} gives the displayed image formula, completing
the proof of all the assertions.
\end{proof}

Removing the common target shift \(D_{q-1}\) presents the matrix as
\[
 R(-e_q)^{r_q}\longrightarrow R^{r_{q-1}}.
\]
Its degree-\(h\) image is the degree-\(h+D_{q-1}\) image in
\eqref{alg:limit-hilbert}. Thus its image dimension is constant for
every \(h\). Taking the coefficient closure of the tuples in each of
finitely many charts, projecting onto stage \(q\), and taking their
finite union gives a compact matrix family with exactly this property.

\begin{remark}
Keeping the scheme structure is essential in this compactness
argument. For example, the sections \(X(X+\varepsilon Y)\) of
\(\mathcal O_{\mathbb P^1}(2)\), after harmless scalar normalization,
can tend to \(X^2\). The limiting doubled point has the same quotient
Hilbert polynomial as the two preceding points. Replacing the limit
by its reduction would change that polynomial and destroy the
dimension argument.
\end{remark}

\section{Step 2: Uniform holomorphic division}
\label{sec:division}

The preceding construction has converted rescaled algebraic data into
a compact family with a fixed Hilbert function. We now show that these
two properties imply a uniform division estimate between two balls.
The argument uses a finite Gr\"obner basis only at one reference
matrix, followed by a bounded perturbation on a coefficient space
containing every degree at once. The leading-term rescaling is related
to the stable division argument of Shalit
\cite[Proposition~3.2, Lemmas~3.4--3.5, and Theorem~3.6]{ShalitDivision}.
Here the parameter uniformity and the treatment of the homogenizing
variable must both be retained.

In this section \(R=\mathbb C[Z_0,\ldots,Z_N]\). Write \(B_s=\{z\in\mathbb C^N:|z|<s\}\), with Lebesgue measure.
For a homogeneous matrix \(A\), put \(A_{\mathrm{aff}}(z)=A(1,z)\).
A holomorphic vector \(F\) belongs to its analytic image sheaf on
\(B_R\) if, near every point of that ball, it has a holomorphic
representation \(F=A_{\mathrm{aff}}c\). This is stronger than
membership at the center, or membership in the image of the evaluated
matrix at each point.

\begin{thmC}[Uniform holomorphic division between two balls]
\label{div:uniform}
Let \(\mathcal A\) be a compact family, in coefficient topology, of
homogeneous matrices
\[
 A:R(-e)^a\longrightarrow R^b,
 \qquad e\geq0,
\]
with fixed sizes and degree. Suppose that for every integer \(h\),
\begin{equation}
 \dim(\operatorname{im}A)_h
 \quad\text{is independent of }A\in\mathcal A.
 \label{div:fixed-function}
\end{equation}
For every \(s>0\), there are \(R>s\) and \(C<\infty\) such that, for
all \(A\in\mathcal A\), every square-integrable holomorphic \(F\) in
the analytic image sheaf on \(B_R\) has a holomorphic solution
\[
 A_{\mathrm{aff}}u=F|_{B_s},
 \qquad
 \|u\|_{L^2(B_s)}\leq C\|F\|_{L^2(B_R)}.
 \tag{\(\dagger\)}
\]
For each fixed $A$, the solution can be chosen linearly in $F$.
In fact the same affine conclusion holds if \eqref{div:fixed-function}
is weakened to constancy of the eventual Hilbert polynomial of
$\operatorname{im}A$; in this weaker formulation no uniform threshold
for equality of the Hilbert functions is assumed.
\end{thmC}

A nonzero homogeneous submodule of a finite free polynomial module
contains a shifted free rank-one submodule, and therefore has nonzero
Hilbert polynomial. Thus if the common polynomial is zero, every
matrix is zero and the theorem is immediate. We prove it otherwise,
for $N\ge1$. In dimension zero the homogeneous image dimensions
in large degrees equal the rank of a constant matrix after setting
$Z_0=1$. Compactness and constant rank bound its positive singular
values away from zero, which gives the assertion directly.

\paragraph{Relation to parameter-dependent division and scissions.}
The proof below is a homogeneous adaptation of the classical bounded
splitting construction of \cite{Galligo,HauserMueller}. In particular,
one may view the perturbation of a reference splitting $\sigma_0$ as
inverting
\[
 U_A=A\sigma_0+(1-A_0\sigma_0)=1+(A-A_0)\sigma_0.
\]
On each degree block $U_A$ carries the reference image injectively into
the new image. Equal dimensions then remove the remainder. We give the
monomial-reconstruction version in detail to verify both the relevant
norms and preservation of homogeneous degree. The selected inverse
need only be complex-linear, not $\mathcal O$-linear. The classical
small-polydisc formulation is compatible with the present large affine
radii: if $r_j$ are homogeneous coefficient radii, write
$\rho_j=r_j/r_0$. For degree-$h$ data and degree-$(h-e)$ coefficients,
\[
 \|F_h\|_{r,\eta}=r_0^h\|F_h\|_{\rho,\eta},\qquad
 \|u_{h-e}\|_r=r_0^{h-e}\|u_{h-e}\|_\rho.
\]
A splitting estimate at fixed $r$ therefore changes only by a factor
$r_0^e$, independent of $h$. The explicit construction below chooses
$\rho_0=1$ and the affine radii directly. In contrast to
Beatrous' restriction theorem \cite{Beatrous} and its extension in
\cite{DouglasTangYu}, the datum here is already an ambient holomorphic
section in the image sheaf. The loss of domain between $B_R$ and $B_s$
is part of the estimate.

\subsection{A reference initial module and coefficient weights}

Fix \(A_0\in\mathcal A\). Order module monomials \(Z^\alpha e_\nu\)
by lexicographic comparison in the variable order
\[
 Z_1,\ldots,Z_N,Z_0,
\]
using a fixed component order to break ties. Choose homogeneous
vectors
\[
 H_i^0=A_0V_i,\qquad 1\leq i\leq k,
\]
with monic leading module monomials \(m_i\) that generate the initial
module of \(\operatorname{im}A_0\). Such a finite collection exists
because monomial submodules of a finite free polynomial module are
finitely generated. The source vector \(V_i\) can be chosen
homogeneous of degree \(\deg m_i-e\), by taking that component of
any polynomial representation of \(H_i^0\).

At each degree \(h\), the module monomials not in this initial module
form a basis of the quotient by \((\operatorname{im}A_0)_h\).
Indeed, division by the \(H_i^0\) terminates in that finite-dimensional
homogeneous block and leaves a combination of those complementary
monomials. Such a nonzero combination cannot lie in the image,
because its leader would then belong to the initial module.
Consequently the number of degree-\(h\) monomials in the initial
module equals \(\dim(\operatorname{im}A_0)_h\).

For positive numbers \(\rho_j\) and \(\eta_\nu\), define weighted
coefficient norms by
\[
 \|f\|_{\rho,\eta}
 =
 \sum_{\alpha,\nu}|f_{\alpha,\nu}|\rho^\alpha\eta_\nu,
 \qquad
 \rho^\alpha=\prod_{j=0}^{N}\rho_j^{\alpha_j}.
 \]
Scalar series have the norm \(\|\cdot\|_\rho\). Source-vector norms
are sums of the scalar norms of their components.

\begin{lemma}[Dominance of the reference leaders]
\label{div:weights}
For the prescribed \(s>0\), there are weights with
\[
 \rho_0=1,\qquad \rho_j>s\quad(1\leq j\leq N),
 \qquad \eta_\nu>0,
\]
and a coefficient neighborhood \(U\) of \(A_0\), such that
\[
 \|H_i^0-m_i\|_{\rho,\eta}\leq\tfrac14\|m_i\|_{\rho,\eta},
 \qquad
 \|AV_i-H_i^0\|_{\rho,\eta}\leq\tfrac14\|m_i\|_{\rho,\eta}
 \tag{\(\ddagger\)}
\]
for every \(i\) and every \(A\in U\).
\end{lemma}

\begin{proof}
Assign integer scores \(c_\nu\) increasing with component priority.
Set \(w_0=0\), and choose positive integers
\(w_N,w_{N-1},\ldots,w_1\) in that order. Only the finitely many
comparisons between a leader of an \(H_i^0\) and a term in its tail
need to be considered. For one whose first different affine exponent
occurs at \(Z_j\), say leader \((\alpha,\nu)\) and tail
\((\beta,\mu)\), require
\[
 w_j>
 \left|\sum_{\ell=j+1}^{N}w_\ell(\alpha_\ell-\beta_\ell)
       +c_\nu-c_\mu\right|.
\]
This is possible at each step. The leader has a larger exponent of
\(Z_j\), so its weighted exponent is strictly larger. If all affine
exponents agree, homogeneity forces equality of the \(Z_0\)-exponents,
and the component scores decide the comparison.

Put \(\rho_j=T^{w_j}\) and \(\eta_\nu=T^{c_\nu}\).
Each tail-to-leader monomial weight ratio is at most \(T^{-1}\).
For sufficiently large \(T\), the sum of its finitely many
coefficient-weighted ratios is at most \(1/4\) for each \(H_i^0\),
and every affine radius exceeds \(s\). The first inequality follows.
With these weights and the \(V_i\) fixed, continuity of
\(A\mapsto AV_i\) in a finite-dimensional coefficient space supplies
one neighborhood \(U\) for all the second inequalities.
\end{proof}

The leading monomial of \(AV_i\) is allowed to change in this
neighborhood. The estimates use the old monomial partition and
control its reconstruction map by a norm perturbation.

\subsection{One bounded inverse for all homogeneous degrees}

Let \(\mathcal W_{\rho,\eta}\) be the Banach space of target formal
series with finite weighted coefficient norm. For a scalar series
\(v\) and target series \(f\), absolute convolution gives
\[
 \|vf\|_{\rho,\eta}\leq\|v\|_\rho\|f\|_{\rho,\eta}.
 \tag{\(\star\)}
\]
Multiplication by a single module monomial gives equality with its
monomial weight.

Partition all target module monomials as follows. A monomial lies
in \(\Delta_i\) if \(m_i\) is the first leader in the chosen list
that divides it; it lies in \(\Delta\) if none does. Divisibility
requires the same target component. Let \(\mathcal Q_i\) consist of
scalar series \(v_i\) for which \(v_im_i\) is supported in
\(\Delta_i\), and let \(\mathcal R_\Delta\) consist of target series
supported in \(\Delta\). These are closed coefficient subspaces.
Give
\[
 \mathcal D=
 \left(\bigoplus_{i=1}^k\mathcal Q_i\right)\oplus\mathcal R_\Delta
\]
the norm
\[
 \|(v,r)\|_{\mathcal D}
 =
 \sum_i\|m_i\|_{\rho,\eta}\|v_i\|_\rho+\|r\|_{\rho,\eta}.
\]
The map
\[
 T_{\mathrm{mon}}(v,r)=\sum_i v_im_i+r
\]
is an isometric bijection onto \(\mathcal W_{\rho,\eta}\), since
its summands occupy disjoint monomial sets.

\begin{lemma}[A uniform, degree-preserving inverse]
\label{div:neumann}
For \(A\in U\), the map
\[
 T_A(v,r)=\sum_i v_iAV_i+r
\]
is invertible and satisfies
\[
 \|T_A^{-1}\|_{\mathcal W_{\rho,\eta}\to\mathcal D}\leq2.
 \]
For homogeneous input of degree \(h\), its inverse has \(r\)
homogeneous of degree \(h\), and \(v_i\) homogeneous of degree
\(h-\deg m_i\).
\end{lemma}

\begin{proof}
Lemma~\ref{div:weights} gives
\(\|AV_i-m_i\|_{\rho,\eta}\leq\frac12\|m_i\|_{\rho,\eta}\).
By \((\star)\),
\[
 \|(T_A-T_{\mathrm{mon}})(v,r)\|_{\rho,\eta}
 \leq\tfrac12\|(v,r)\|_{\mathcal D}.
\]
Thus
\[
 K_A=(T_A-T_{\mathrm{mon}})T_{\mathrm{mon}}^{-1}
 \quad\text{satisfies}\quad \|K_A\|\leq\tfrac12,
\]
and
\[
 T_A^{-1}
 =
 T_{\mathrm{mon}}^{-1}\sum_{\ell=0}^{\infty}(-K_A)^\ell.
\]
The series converges in operator norm and has norm at most two.

Every \(AV_i\) is homogeneous of degree \(\deg m_i\).
Consequently \(T_A\) intertwines the projection onto total degree
\(h\) in the target with the projection onto degrees
\(h-\deg m_i\) and \(h\) in its domain. Its inverse intertwines the
same projections. Those homogeneous blocks are finite dimensional,
so the asserted components are polynomials.
\end{proof}

\begin{lemma}[Vanishing of the remainder on the image]
\label{div:polynomial}
For \(A\in U\cap\mathcal A\) and any degree \(h\) in which
\(\dim(\operatorname{im}A)_h=\dim(\operatorname{im}A_0)_h\), every
homogeneous \(F\in(\operatorname{im}A)_h\) admits a homogeneous source vector
\(u\) of degree \(h-e\) satisfying
\begin{equation}
 Au=F,\qquad
 \|u\|_\rho\leq C_0\|F\|_{\rho,\eta},
 \qquad
 C_0=2\max_i\frac{\|V_i\|_\rho}{\|m_i\|_{\rho,\eta}}.
 \label{div:polynomial-bound}
\end{equation}
The same \(C_0\) works in every degree satisfying this equality; under
\eqref{div:fixed-function} it works in all degrees.
\end{lemma}

\begin{proof}
Let \(\mathcal Q_h\) be the degree-\(h\) part of
\(\bigoplus_i\mathcal Q_i\), where \(v_i\) is assigned shifted
degree \(\deg v_i+\deg m_i\).
Its dimension is the number of degree-\(h\) monomials in the
reference initial module, namely
\(\dim(\operatorname{im}A_0)_h\).
The map \(T_A\) sends \((\mathcal Q_h,0)\) injectively into
\((\operatorname{im}A)_h\), since each \(AV_i\) is a column
combination of \(A\). The stated dimension equality gives
\[
 \dim\mathcal Q_h
 =
 \dim(\operatorname{im}A_0)_h
 =
 \dim(\operatorname{im}A)_h.
\]
The image is therefore the whole target image. The unique inverse
of \(F\) under \(T_A\) has the form \((v,0)\).

Set \(u=\sum_i v_iV_i\). Its summands all have degree \(h-e\), and
\[
 \|u\|_\rho
 \leq\sum_i\|v_i\|_\rho\|V_i\|_\rho
 \leq
 \left(\max_i\frac{\|V_i\|_\rho}{\|m_i\|_{\rho,\eta}}\right)
 \|(v,0)\|_{\mathcal D}.
\]
Lemma~\ref{div:neumann} proves the asserted bound.
\end{proof}

The all-degree hypothesis gives this estimate in every degree. It is
not equivalent to constancy of the eventual Hilbert polynomial. For
example, $A_t=(X^2,tXY,Y^2)$ on $\mathbb P^1$, $|t|\le1$, has the
same image Hilbert polynomial for all $t$: in degrees $h\ge3$ the
image is the entire homogeneous space. But in degree two the image
has dimension three for $t\ne0$ and two for $t=0$. Representing $XY$
in degree two requires coefficient $1/t$ on the middle entry. No
uniform all-degree estimate in the fixed coefficient norm is possible.
Over $\C\{t\}$ the quotient by $(X^2,tXY,Y^2)$ has the nonzero
$t$-torsion class of $XY$ and is not flat. This explains why one
cannot simply invoke Galligo's analytic flatness criterion from
constancy of an eventual projective Hilbert polynomial. The affine
conclusion of C nevertheless holds, as the next argument shows.

\subsection{Dehomogenization without a degree-dependent loss}

Write \(\rho'=(\rho_1,\ldots,\rho_N)\). For an affine polynomial
vector \(f(z)=\sum_\beta f_\beta z^\beta\) and \(h\geq\deg f\),
define its degree-\(h\) homogenization by
\[
 f^{[h]}(Z_0,Z')
 =
 \sum_\beta f_\beta Z_0^{h-|\beta|}(Z')^\beta.
\]
Because \(\rho_0=1\),
\begin{equation}
 \|f^{[h]}\|_{\rho,\eta}=\|f\|_{\rho',\eta}.
 \label{div:homogenization-norm}
\end{equation}
There is exactly one \(Z_0\)-power for each affine monomial in a
fixed homogeneous degree, so dehomogenization causes no merging of
distinct coefficients.

\begin{lemma}[Affine polynomial division]
\label{div:affine}
For \(A\in U\cap\mathcal A\), every affine polynomial vector
\(F\in A_{\mathrm{aff}}\mathbb C[z]^a\) has a polynomial preimage
\(u\) satisfying
\[
 A_{\mathrm{aff}}u=F,\qquad
 \|u\|_{\rho'}\leq C_0\|F\|_{\rho',\eta}.
 \]
\end{lemma}

\begin{proof}
Choose a polynomial representation \(F=A_{\mathrm{aff}}c\) and
an integer \(h\geq e+\deg c\). If only the eventual Hilbert
polynomial is fixed, increase \(h\) so that the image dimensions of
\(A\) and \(A_0\) agree in degree \(h\); such a degree exists for this
particular \(A\). No common threshold is needed. Both sides of
\[
 F^{[h]}=A\,c^{[h-e]}
\]
are homogeneous of degree \(h\) and dehomogenize to the same
polynomial, so the equality holds. Apply
Lemma~\ref{div:polynomial} and dehomogenize its preimage. Equation
\eqref{div:homogenization-norm}, also with unweighted source
components, gives the bound. The choice of \(h\) introduces no
factor into the estimate.
\end{proof}

The representation degree can be much larger than the degree of
the datum. For example, the row \((z,z+1)\) maps \((-1,1)^T\) to
the constant \(1\), while adding any multiple of
\((z+1,-z)^T\) produces representations of arbitrarily high degree.
The unit weight on the homogenizing variable is what makes the
estimate insensitive to this choice.

\subsection{Holomorphic data and the two-ball estimate}

For a positive vector of radii \(\tau\), put
\[
 P_\tau=\{z\in\mathbb C^N:|z_j|<\tau_j\text{ for every }j\}.
\]
For the reference neighborhood \(U\), set
\[
 R_U=2|\rho'|_2+1.
\]
The domains satisfy
\[
 B_s\subset P_{\rho'}\Subset P_{2\rho'}\Subset B_{R_U}.
\]

\begin{lemma}[Holomorphic division in the coefficient norm]
\label{div:holomorphic}
For \(A\in U\cap\mathcal A\), every holomorphic section \(F\) of
the analytic image sheaf on \(B_{R_U}\) has a holomorphic preimage
\(u\) on \(P_{\rho'}\) with
\[
 A_{\mathrm{aff}}u=F,\qquad
 \|u\|_{\rho'}\leq C_0\|F\|_{\rho',\eta}.
 \]
Here the norms are sums of the Taylor coefficients at the origin.
\end{lemma}

\begin{proof}
On \(B_{R_U}\), let \(\mathscr J_A\) be the analytic image and
\(\mathscr Z_A\) the kernel of the holomorphic matrix. Oka
coherence makes \(\mathscr Z_A\) coherent, and
\[
 0\longrightarrow\mathscr Z_A\longrightarrow\mathcal O^a
 \xrightarrow{A_{\mathrm{aff}}}\mathscr J_A
 \longrightarrow0
\]
is exact. The ball is Stein, so Cartan's theorem B gives
\(H^1(B_{R_U},\mathscr Z_A)=0\). The global-section map is
therefore onto: some holomorphic \(c\) on \(B_{R_U}\) satisfies
\(A_{\mathrm{aff}}c=F\). These are the standard coherence and
Stein vanishing theorems; see \cite{GrauertRemmert}.

Let \(c^{(k)}\) be the Taylor polynomial of \(c\) through total
degree \(k\). The closure of \(P_{2\rho'}\) is compact in the
input ball. Cauchy's estimate bounds each component coefficient
by a constant times \((2\rho')^{-\beta}\). Since
\[
 \sum_{\beta\in\mathbb N^N}2^{-|\beta|}=2^N,
\]
the sequence \(c^{(k)}\) converges to \(c\) in
\(\|\cdot\|_{\rho'}\). Multiplication by the fixed polynomial
matrix is continuous in the coefficient norms; hence
\[
 F^{(k)}=A_{\mathrm{aff}}c^{(k)}
 \longrightarrow F
 \quad\text{in }\|\cdot\|_{\rho',\eta}.
\]
These approximants lie in the polynomial image.

Apply Lemma~\ref{div:affine} separately to each \(F^{(k)}\),
obtaining polynomials \(u^{(k)}\) with
\[
 A_{\mathrm{aff}}u^{(k)}=F^{(k)},\qquad
 \|u^{(k)}\|_{\rho'}\leq C_0\|F^{(k)}\|_{\rho',\eta}.
\]
Every coefficient is bounded. A diagonal subsequence converges
coefficientwise to a series \(u\), and Fatou's lemma for the
weighted sums gives
\[
 \|u\|_{\rho'}\leq C_0\|F\|_{\rho',\eta}.
\]
For \(0<\theta<1\), the tail of total degree greater than \(d\)
on \(P_{\theta\rho'}\) is bounded by
\[
 \theta^{d+1}\sup_k\|u^{(k)}\|_{\rho'}.
\]
The same holds for the limit. Finite coefficient convergence
therefore gives locally uniform convergence on \(P_{\rho'}\).
Its limit is holomorphic and satisfies the required equation.

No quantitative bound on the initial Cartan-B lift \(c\) was
used. It supplied polynomial image approximants; the bounded
lifts were then constructed afresh from their data.
\end{proof}

Approximating the datum itself by Taylor polynomials would not be
valid here. For example, if
\[
 A_{\mathrm{aff}}(z)=z-1,\qquad F(z)=(z-1)e^z,
\]
then the total Taylor polynomial \(T_kF\) has value \(-1/k!\) at
\(z=1\), and so is not in the polynomial image. The approximants
\((z-1)T_k(e^z)\) used by the proof do belong to that image.

\begin{lemma}[Conversion to ball norms]
\label{div:ball}
For every \(A\in U\cap\mathcal A\), the conclusion of
Theorem~\ref{div:uniform} holds with input radius \(R_U\) and
constant
\[
 C_U=
 C_0\,2^N\left(\sum_{\nu=1}^b\eta_\nu\right)
 \left(\frac{\operatorname{Vol}(B_s)}
 {\operatorname{Vol}(B_{1/2})}\right)^{1/2}.
\]
\end{lemma}

\begin{proof}
Cauchy's estimate on \(P_{2\rho'}\) gives
\[
 \|F\|_{\rho',\eta}
 \leq
 2^N\left(\sum_{\nu=1}^b\eta_\nu\right)
 \sup_{P_{2\rho'}}|F|.
\]
For every point in the closure of this polydisc, its ball of
radius \(1/2\) is contained in \(B_{R_U}\). The submean inequality
for the subharmonic function \(|F|^2\) bounds the supremum by
\[
 \operatorname{Vol}(B_{1/2})^{-1/2}\|F\|_{L^2(B_{R_U})}.
\]
On \(B_s\subset P_{\rho'}\), the Euclidean norm of the preimage
from Lemma~\ref{div:holomorphic} is at most its source
coefficient norm. Hence
\[
 \|u\|_{L^2(B_s)}
 \leq\operatorname{Vol}(B_s)^{1/2}\|u\|_{\rho'}.
\]
Combining the three inequalities gives the claim.
\end{proof}

\begin{proof}[Completion of the proof of Theorem~\ref{div:uniform}]
The neighborhoods \(U\) obtained from all reference matrices
cover the compact family \(\mathcal A\). Choose a finite
subcover \(U_1,\ldots,U_M\), and take
\[
 R=\max_j R_{U_j},\qquad C=\max_j C_{U_j}.
\]
For \(A\in U_j\), restrict an admissible datum from \(B_R\) to
\(B_{R_{U_j}}\). Analytic image membership is preserved, and
its norm can only decrease. Lemma~\ref{div:ball} supplies the
solution on the same prescribed output ball \(B_s\).

For a fixed \(A\), multiplication by \(A_{\mathrm{aff}}\) is a
bounded map between the appropriate vector-valued Bergman
spaces on \(B_s\), since its entries are polynomials. Its
kernel is closed. The nonempty fiber over \(F|_{B_s}\) is an
affine translate of this kernel; it has a unique element
perpendicular to the kernel. This element minimizes the norm,
so it obeys the same estimate. Sums and scalar multiples of
these perpendicular solutions solve the corresponding equations
and remain perpendicular to the fixed kernel. Uniqueness proves
linearity on the space of admissible data on \(B_R\).
Closedness of the range is neither asserted nor needed.
\end{proof}

\begin{remark}[What the theorem requires]
The theorem supplies a sufficiently large input radius for a
prescribed output radius. It does not assert the estimate for
every prescribed pair of radii. It also requires analytic image
membership throughout the input ball. For instance, \(1\) is
divisible by \(z-t\) near the origin when \(t\neq0\), but has no
holomorphic quotient on a ball containing \(t\). Finally,
pointwise membership is insufficient: \(z\) belongs to the
image of the scalar matrix \(z^2\) at each point, including
zero, but has no holomorphic quotient near zero.

The compactness hypothesis alone, even with independent limiting columns,
does not ensure uniform division. Consider the homogeneous rows
$(X^2,XY+tY^2)$, $|t|\leq1$, on $\Pj^1$. On $Y=1$ the row is
$(z^2,z+t)$. For $t\ne0$, the constant $1$ has a polynomial
representation, for instance
$1=t^{-2}(z^2-(z-t)(z+t))$. If $0<|t|<s/2$, every holomorphic
solution of $z^2u+(z+t)v=1$ on $B_s$ satisfies $u(-t)=t^{-2}$.
The submean inequality on the disk of radius $s/2$ centered at $-t$
gives $\|u\|_{L^2(B_s)}\geq\sqrt\pi(s/2)|t|^{-2}$ for ordinary
area measure. No uniform two-ball constant exists as $t\to0$.
Indeed the projective image is the unit ideal sheaf for $t\ne0$,
but at $t=0$ it is the ideal sheaf of a point. Its Hilbert polynomial
changes, precisely the kind of degeneration excluded by
Theorem~\ref{thm:B} for the normalized families used here.
\end{remark}

For the normalized presentations, Theorem~\ref{alg:hilbert}
verifies the fixed-Hilbert-function hypothesis, and
Proposition~\ref{alg:compact} supplies compactness. Each original
or limiting presentation is exact on holomorphic stalks at
$E_0,\ldots,E_{N-1}$, and its first row generates its corresponding
ideal sheaf. No injectivity of the last map is asserted. Consequently
both ideal data and the holomorphic cycles in these stages admit the
uniform local division just proved. Proposition~\ref{ana:microscopic-division}
uses exactly these stalks; the final correction form has degree $(0,N)$
and is automatically $\bar\partial$-closed. The next step is to
transport this estimate back from the frozen Euclidean
coefficient norms to metrics on the actual bundles.

% END sections/algebra.tex


% BEGIN sections/analysis.tex
\section{Metrics adapted to the presentation}
\label{ana:metrics}

We now turn the algebraic compactness of the preceding section into
estimates in the original scalar line norm.  Throughout this section
$N\geq1$, $W=\Pj^N$, and $L=\cO_W(1)$ has its Fubini--Study metric.
Fix the complete presentation of Proposition~\ref{alg:complete}, through
stage $N$.  Thus $E_{-1}=\cO_W$, $E_q=V_q\otimes L^{-D_q}$,
$d_0=\Phi$, and $d_q:E_q\to E_{q-1}$.  Thus $\im d_0=\cI$ and
$\im d_q=\ker d_{q-1}$ for $1\leq q\leq N$ on holomorphic stalks;
there is no exactness requirement at the last source $E_N$.
The bundles, their bases, their degrees, and the maps are fixed before any
analytic parameter is chosen.  Zero coefficient spaces are allowed, and
all forms and maps at a zero stage are understood to be zero.

Choose a finite collection of buffered affine charts
$V_\alpha\Subset U_\alpha$, with the $V_\alpha$ covering $W$.
We use standard projective line frames on these charts, so the matrices
$d_{q,\alpha}$ are polynomials of fixed degrees.  Every compact set on
which a chart is used has a fixed positive buffer inside that chart.
After decreasing $\varepsilon_0>0$, all sampling balls of radius
$0<\varepsilon\leq\varepsilon_0$ centered in these sets remain in their
charts.  Further fixed multiples of this radius will be accommodated
by decreasing $\varepsilon_0$ when they are introduced.

Recall the local Gram matrices of Lemma~\ref{alg:gram}:
\begin{equation}\label{ana:gram-recall}
 G_{-1,\alpha,\varepsilon}=1,\qquad
 G_{q,\alpha,\varepsilon}(p)
 =\int_{B_1}d_{q,\alpha}(p+\varepsilon z)^*
 G_{q-1,\alpha,\varepsilon}(p)
 d_{q,\alpha}(p+\varepsilon z)\,d\mu(z),
\end{equation}
where $\mu$ is probability volume on the unit ball in $\C^N$.
Equivalently, if
\[
 P_q(p;z_0,\ldots,z_q)
 =d_{0,\alpha}(p+\varepsilon z_0)\cdots
 d_{q,\alpha}(p+\varepsilon z_q),
\]
then
\begin{equation}\label{ana:product-gram}
 G_{q,\alpha,\varepsilon}(p)[v]
 =\int_{(B_1)^{q+1}}|P_q(p;z_0,\ldots,z_q)v|^2\,d\mu^{q+1}.
\end{equation}
The sample points are independent. For $q\geq1$, evaluating every map
at the same point would make the products vanish by the complex identities.
The constant independence of the columns of each complete evaluation
instead makes \eqref{ana:gram-recall} positive definite at every
positive radius, as proved in Lemma~\ref{alg:gram}.

\begin{thmD}[Compatible metrics with controlled curvature]\label{thm:D}
For the presentation in Theorem~\ref{thm:B} there are smooth positive
Hermitian metrics $H_{q,\varepsilon}$ on the nonzero bundles $E_q$, with
$H_{-1,\varepsilon}=1$, for all sufficiently small $\varepsilon>0$.
On each smaller buffered chart $V_\alpha\Subset U_\alpha$,
they are uniformly comparable on the actual fibers to the sampled
Gram metrics $G_{q,\alpha,\varepsilon}$: a common constant $C\geq1$
satisfies
\[
 C^{-1}G_{q,\alpha,\varepsilon}(p)[v]
 \leq H_{q,\varepsilon}(p)[v]
 \leq C G_{q,\alpha,\varepsilon}(p)[v]
 \qquad(p\in V_\alpha,\ v\in(E_q)_p).
\]
No such uniform comparison on the full affine charts is claimed.
Each map $d_q$ is contractive from its source metric to its target metric.
With curvature measured in the indicated fiber and Fubini--Study tangent metrics,
\begin{align}
 |\Phi v|&\le |v|_{H_{0,\varepsilon}},\label{eq:geo-contract}\\
 |(\partial_{z_j}\Phi)v|_{L^{-1}}
 &\le C\varepsilon^{-1}|v|_{H_{0,\varepsilon}},\label{eq:geo-differentiate}\\
 \|\Theta(H_{q,\varepsilon})\|_{H_{q,\varepsilon},\omega_{\FS}}
 &\le C\varepsilon^{-2}.\label{eq:geo-curvature}
\end{align}
The constants depend on the fixed presentation and geometry, not on the
scale or the point. In particular the absolute value of each Nakano form
is bounded by a constant times $\varepsilon^{-2}$ times the squared
fiber norm of its arguments.
\end{thmD}

\paragraph{Why the metric varies.}
Take the presentation $Z_0:\cO_{\mathbb P^1}(-1)\to\cO$ of the point
$Z_0=0$. For $F_n=Z_0Z_1^n$ the holomorphic lift is unique and equals
$a_n=Z_1^n$. In the standard Fubini--Study metrics,
\[
 \|a_n\|^2=\frac1{n+1},\qquad
 \|F_n\|^2=\frac1{(n+1)(n+2)}.
\]
Consequently no fixed smooth positive metric on the coefficient line can
make the division estimate uniform: all such metrics are comparable to
the standard one on compact projective space. But the metric
\[
 |v|_{H_\varepsilon}^2=|Z_0v|^2+\varepsilon^2|v|_{\rm std}^2
\]
gives, after twisting by $L^{n+1}$,
\[
 \|a_n\|_{H_\varepsilon\otimes L^{n+1}}^2
 = (1+\varepsilon^2(n+2))\|F_n\|^2.
\]
The scale $\varepsilon\asymp n^{-1/2}$ repairs the division estimate.
Such a change alone is not enough: it must also control derivatives of
the map out of the coefficient bundle and curvature of the new metric.
This is the reason for the particular sampling construction below.

\paragraph{Comparison with singularly weighted division.}
Skoda's $L^2$ theorem \cite[Theorem~1]{SkodaDivision}, for a row $g$ of
holomorphic functions on a pseudoconvex domain in $\C^N$, uses on its
right side the weight $|g|^{-2(\alpha q+1)}e^{-\psi}$, where
$q=\min(N,r-1)$, $r$ is the number of generators, $\alpha>1$, and $\psi$
is plurisubharmonic. Ordinary ideal membership need not imply finiteness
of this integral. For example, $g=(z_1^2,z_2^2)$, $F=z_1^2$ and $\psi=0$
in $\C^2$ give, on a set of angular directions of positive measure, the
radial integral $\int_0^1 t^{3-4\alpha}\,dt=\infty$. The desired
right side in Theorem E is instead the original scalar norm of $F$.
This does not rule out other approaches using refined $L^2$ division
\cite{SkodaBundles,Varolin,Ji}; it identifies the specific norm problem
that the present metrics address.

\paragraph{Sampling as a finite-jet metric.}
On one affine chart, expand the polynomial row about $p$ and use the
orthogonality of monomials on the probability unit ball. The exact identity
is
\begin{equation}\label{eq:sampled-jets}
 G_{0,\varepsilon}(p)[v]=
 \sum_{|\alpha|\le D_0}
 \frac{N!\,\varepsilon^{2|\alpha|}}{\alpha!(N+|\alpha|)!}
 |(\partial^\alpha\Phi)(p)v|^2.
\end{equation}
Here $\Phi$ denotes the dehomogenized row in that chart. Indeed the
Taylor coefficient of $z^\alpha$ is $\varepsilon^{|\alpha|}
(\partial^\alpha\Phi)(p)v/\alpha!$, and the unit-ball moment is
$N!\alpha!/(N+|\alpha|)!$. Thus the metric sees all the finite jets at
the sampling scale, even where the value of the row vanishes. The
orders zero and one immediately control evaluation and first derivative
evaluation. At higher stages one must use independent sample points,
since same-point sampling of consecutive differentials vanishes.
For instance the two maps $(x,y)$ and $(-y,x)^T$ on $\C^2$ give
$G_{0,\varepsilon}(0)=\varepsilon^2 I_2/3$ and
$G_{1,\varepsilon}(0)=2\varepsilon^4/9$, by the second moments
of independent points in the unit ball. These small eigenvalues are
not an obstruction: the derivative estimates will be measured relative
to the metrics themselves.

\subsection{Polynomial estimates and relative derivatives}

\begin{lemma}\label{ana:polynomial-estimates}
Fix a bound for the degree and two positive radii.  The $L^2$ norms on
the corresponding balls are equivalent on the space of polynomials of
that degree.  Evaluation at a fixed point in a fixed bounded set,
differentiation as a map between polynomial spaces, and derivative
evaluation are bounded in these norms.  The constants remain valid
for polynomials with values in any finite-dimensional Hilbert space,
independently of its inner product.  The same assertions hold on
products of balls with fixed degrees in each group of variables.
\end{lemma}

\begin{proof}
Each ball integral is a norm on one finite-dimensional scalar polynomial
space: vanishing of the integral forces the polynomial to vanish on an
open set and hence identically.  Norm equivalence and boundedness of
linear maps between finite-dimensional spaces give the scalar claims.
For a vector-valued polynomial choose an orthonormal basis of the given
target Hilbert space, apply the scalar estimate to every component,
and sum the squared inequalities.  This procedure introduces no
constant depending on the target metric or its smallest eigenvalue.
Apply the same argument successively in the different variable groups
to obtain the assertions on products.
\end{proof}

\begin{lemma}\label{ana:relative-gram}
In a fixed chart put $G(p)=G_{q,\alpha,\varepsilon}(p)$ and
$|v|_p^2=v^*G(p)v$.  For constant coefficient vectors $v,w$ and affine
coordinate indices $a,b$,
\begin{align}
 |v^*(\partial_aG)w|&\leq C\varepsilon^{-1}|v|_p|w|_p,
 \label{ana:first-relative}\\
 |v^*(\partial_{\bar b}\partial_aG)w|
 &\leq C\varepsilon^{-2}|v|_p|w|_p.
 \label{ana:second-relative}
\end{align}
If two centers in a buffered chart are joined by a segment of length
at most $A\varepsilon$, their Gram norms, expressed in that frame,
are comparable with constants depending only on $A$ and the fixed
presentation.  Furthermore,
\begin{align}
 |d_q(p)v|_{G_{q-1}(p)}&\leq C|v|_{G_q(p)},
 \label{ana:local-map}\\
 |(\partial_a d_q)(p)v|_{G_{q-1}(p)}
 &\leq C\varepsilon^{-1}|v|_{G_q(p)}.
 \label{ana:local-map-derivative}
\end{align}
At $q=0$ the frozen target norm is the usual absolute value in the
local frame of the trivial line bundle.
\end{lemma}

\begin{proof}
Let $F(p)$ send $v$ to the product polynomial $P_q(p;\cdot)v$ in
the fixed Hilbert space of product polynomials with its probability
$L^2$ norm.  Then $G(p)=F(p)^*F(p)$ and $F$ depends holomorphically
on $p$.  Each appearance of $p_a$ in the $j$th factor is through
$p_a+\varepsilon z_{j,a}$.  Consequently
\[
 \partial_{p_a}(F(p)v)
 =\varepsilon^{-1}\sum_{j=0}^q
 \partial_{z_{j,a}}(F(p)v).
\]
The degree bounds and the number of groups are fixed, so
Lemma~\ref{ana:polynomial-estimates} gives
\[
 \|(\partial_aF(p))v\|\leq C\varepsilon^{-1}\|F(p)v\|.
\]
Holomorphic dependence gives the exact matrix identities
\[
 \partial_aG=F^*\partial_aF,
 \qquad
 \partial_{\bar b}\partial_aG=(\partial_bF)^*\partial_aF.
\]
Cauchy--Schwarz proves \eqref{ana:first-relative} and
\eqref{ana:second-relative}.

For a nonzero constant vector $v$, the first estimate and its conjugate
bound the absolute derivative of $\log(v^*G(\gamma(t))v)$ along a
segment by $C|\gamma'(t)|/\varepsilon$.  Integration gives factors
$e^{-CA}$ and $e^{CA}$ for the squared norms.  The buffers ensure
that the segments under consideration stay in the region where the
Gram matrices are defined.

Finally give the target of the polynomial
$z\mapsto d_q(p+\varepsilon z)v$ the frozen norm $G_{q-1}(p)$.
Its squared $L^2(B_1)$ norm is $G_q(p)[v]$ by
\eqref{ana:gram-recall}.  Evaluation at zero proves
\eqref{ana:local-map}.  Derivative evaluation, followed by division
by $\varepsilon$, proves \eqref{ana:local-map-derivative}.
The vector-valued assertion of Lemma~\ref{ana:polynomial-estimates}
is what makes these estimates uniform even when a Gram eigenvalue
tends to zero with $\varepsilon$.
\end{proof}

\subsection{Comparison on actual bundle fibers}

The local quadratic forms must be compared on a vector in $E_q$, not
on two unrelated coordinate columns.  Set $D_{-1}=0$.  On the overlap
of two line frames write
\[
 e_\beta(x)=t(x)e_\alpha(x),
\]
where $t$ is holomorphic and nowhere zero.  If $v\in(E_q)_p$ has
column $v_\alpha$ in the first frame, it has column
$v_\beta=t(p)^{D_q}v_\alpha$ in the second.  The matrix transition is
\begin{equation}\label{ana:map-transition}
 d_{j,\beta}(x)=t(x)^{D_{j-1}-D_j}d_{j,\alpha}(x),
\end{equation}
with the corresponding change of affine coordinates understood.

\begin{proposition}\label{ana:fiber-comparison}
After a common decrease of $\varepsilon_0$, the Gram forms defined
in any two charts satisfy, for $p\in V_\alpha\cap V_\beta$,
\[
 C^{-1}G_{q,\alpha,\varepsilon}(p)[v]
 \leq G_{q,\beta,\varepsilon}(p)[v]
 \leq C G_{q,\alpha,\varepsilon}(p)[v]
\]
for every actual fiber vector $v\in(E_q)_p$.  The constant is
uniform in $p$, $v$, and $0<\varepsilon\leq\varepsilon_0$.
\end{proposition}

\begin{proof}
The closures of the smaller charts are compact inside their larger
charts.  There are finitely many chart pairs.  We may therefore require
that both systems of sample points at a point in a smaller overlap
remain in a common buffered overlap.  The transitions, their inverses,
and the coordinate Jacobians are uniformly bounded there.

At independent sample points $x_0,\ldots,x_q$, formula
\eqref{ana:map-transition} and the transformation of the input vector
show that the two product evaluations differ by the scalar
\begin{equation}\label{ana:independent-transition}
 T(p;x_0,\ldots,x_q)
 =t(p)^{D_q}\prod_{j=0}^q t(x_j)^{D_{j-1}-D_j}.
\end{equation}
All factors are evaluated at their indicated points; no telescoping
at one point is being asserted.  Both $|T|$ and $|T|^{-1}$ have
uniform bounds on these overlaps.

Change each sampling variable in the $\beta$ integral to $\alpha$
coordinates.  Its coordinate ball of radius $\varepsilon$ contains
a ball of radius $c\varepsilon$ about the same center and is contained
in a ball of radius $C\varepsilon$, with uniform Jacobian bounds.
The probability-volume normalizations contribute $\varepsilon^{-2N}$
in each variable.  After bounding \eqref{ana:independent-transition}
and the densities, we have upper and lower bounds by integrals of the
same $\alpha$-coordinate product polynomial over products of those
two balls.  Rescale every variable by $\varepsilon$ about the center.
Lemma~\ref{ana:polynomial-estimates} compares the resulting integrals
on the fixed products of balls with the unit-product integral in
\eqref{ana:product-gram}.  This proves both inequalities.
\end{proof}

\paragraph{Why the buffers are needed.}
Consider again $Z_0:\cO_{\Pj^1}(-1)\to\cO$, now on the two full
standard affine charts $U_0=\{Z_0\ne0\}$ and $U_1=\{Z_1\ne0\}$.
For $p$ in their overlap write $w=Z_0/Z_1$ and take the actual
fiber vector $v=Z_0^{-1}(p)=w^{-1}Z_1^{-1}(p)$. The matrix of
$Z_0$ is $1$ in the $U_0$ frame and $w$ in the $U_1$ frame.
With probability area $\mu$ on the unit disk (the frame labels
$0,1$ below denote the two charts), one has
\[
 \begin{split}
 G_{0,0,\varepsilon}(p)[v]&=1,\\
 G_{0,1,\varepsilon}(p)[v]
 &=|w|^{-2}\int_{|\zeta|<1}|w+\varepsilon\zeta|^2\,d\mu(\zeta)
 =1+\frac{\varepsilon^2}{2|w|^2}.
 \end{split}
\]
For fixed $\varepsilon>0$ the ratio is unbounded as $w\to0$.
Thus uniform comparability to a common global metric on both full
charts would be false. Restriction to buffered charts is essential.

Choose a fixed smooth partition of unity $\{\rho_\alpha\}$,
with $\supp\rho_\alpha\subset V_\alpha$, and define metrics on the
actual bundles by
\begin{equation}\label{ana:patched-metric}
 \widehat G_{q,\varepsilon}
 =\sum_\alpha\rho_\alpha G_{q,\alpha,\varepsilon},
 \qquad \widehat G_{-1,\varepsilon}=1.
\end{equation}
These are smooth positive metrics.  Proposition~\ref{ana:fiber-comparison}
makes the Gram metric of every chart $V_\alpha$ uniformly comparable
to the patched metric on $V_\alpha$: every nonzero summand at $p$
is from a chart $V_\beta$ containing $p$. The same pairwise comparison
applies on $V_\alpha\cap\supp\rho_\beta$, including a point where
$\rho_\beta$ itself is zero, since $\supp\rho_\beta\subset V_\beta$.

\begin{proposition}\label{ana:patched-derivatives}
In any of the fixed smaller affine charts, put
$K=\widehat G_{q,\varepsilon}$.  Then
\begin{align*}
 |v^*(\partial_aK)w|&\leq C\varepsilon^{-1}|v|_K|w|_K,\\
 |v^*(\partial_{\bar b}\partial_aK)w|
 &\leq C\varepsilon^{-2}|v|_K|w|_K.
\end{align*}
With $K_- =\widehat G_{q-1,\varepsilon}$ one also has
\[
 |d_qv|_{K_-}\leq C|v|_K,
 \qquad
 |(\partial_a d_q)v|_{K_-}
 \leq C\varepsilon^{-1}|v|_K.
\]
\end{proposition}

\begin{proof}
In an $\alpha$ frame, a second local metric has matrix
\[
 B_\beta(z)=|t(z)|^{2D_q}
 G_{q,\beta,\varepsilon}(y(z)),
\]
where $y$ is the holomorphic coordinate transition.  Its first and
mixed second derivatives are combinations of the corresponding
derivatives of $G_{q,\beta,\varepsilon}$ and derivatives of the
fixed coordinate and scalar frame factors.  All the latter factors
are bounded on the buffered overlap.  Holomorphicity of $y$ preserves
the mixed derivative type.  Lemma~\ref{ana:relative-gram} therefore
gives the asserted orders, measured initially in $B_\beta$ and then
in $K$ by Proposition~\ref{ana:fiber-comparison}.

The mixed derivative of a weighted summand is explicitly
\[
 (\partial_{\bar b}\partial_a\rho_\beta)B_\beta
 +(\partial_a\rho_\beta)\partial_{\bar b}B_\beta
 +(\partial_{\bar b}\rho_\beta)\partial_aB_\beta
 +\rho_\beta\partial_{\bar b}\partial_aB_\beta.
\]
The derivatives of this fixed partition are bounded.  The support
comparison applies even to the first term when $\rho_\beta$ vanishes.
Since $\varepsilon\leq1$, terms of lower order are bounded by the
same $C\varepsilon^{-2}$.  Summing the finitely many terms proves
the metric estimates.  The map estimates follow from
\eqref{ana:local-map}--\eqref{ana:local-map-derivative} and the same
comparison in source and target.  A differentiation of a line-frame
factor adds a bounded multiple of $d_q$, already covered by the first
map estimate.
\end{proof}

\subsection{Graph metrics and curvature}

Define $H_{-1,\varepsilon}=1$ and, successively on actual fibers,
\begin{equation}\label{ana:graph-definition}
 |v|_{H_{q,\varepsilon}}^2
 =|v|_{\widehat G_{q,\varepsilon}}^2
  +|d_qv|_{H_{q-1,\varepsilon}}^2.
\end{equation}
The identity $d_{q-1}d_q=0$ reduces this recursive expression to
\begin{equation}\label{ana:graph-matrix}
 H_{q,\varepsilon}
 =\widehat G_{q,\varepsilon}
   +d_q^*\widehat G_{q-1,\varepsilon}d_q.
\end{equation}
For $q=0$ the last metric is the scalar metric $1$.

\begin{proposition}\label{ana:graph-curvature}
The metrics $H_{q,\varepsilon}$ are smooth and positive, and
\[
 \widehat G_{q,\varepsilon}\leq H_{q,\varepsilon}
 \leq C\widehat G_{q,\varepsilon}.
\]
Every $d_q:(E_q,H_{q,\varepsilon})\to
(E_{q-1},H_{q-1,\varepsilon})$ is contractive.  The Chern curvature
coefficients, measured in the fiber metric $H_{q,\varepsilon}$ and
unitary Fubini--Study tangent coordinates, have operator norm at most
$B\varepsilon^{-2}$.  In particular their Nakano quadratic form has
absolute value at most
\[
 NB\varepsilon^{-2}\sum_{a=1}^N|\xi_a|_{H_{q,\varepsilon}}^2
\]
on a tuple of fiber vectors $(\xi_1,\ldots,\xi_N)$.
\end{proposition}

\begin{proof}
The comparison follows from \eqref{ana:graph-matrix} and
Proposition~\ref{ana:patched-derivatives}.  Contractivity follows
directly from \eqref{ana:graph-definition}.  It persists after a
common line twist and after tensoring with the same exterior-form
metric.

For the curvature estimate write, in a fixed holomorphic frame,
$H=J+D^*KD$, with $D=d_q$,
$J=\widehat G_{q,\varepsilon}$, and
$K=\widehat G_{q-1,\varepsilon}$.  Use $K=1$ at stage zero.
If subscripts denote chart derivatives, holomorphicity of $D$ gives
\begin{align*}
 (D^*KD)_a&=D^*K_aD+D^*KD_a,\\
 (D^*KD)_{a\bar b}
 &=D_b^*K_aD+D^*K_{a\bar b}D
   +D_b^*KD_a+D^*K_{\bar b}D_a.
\end{align*}
We have $|Dv|_K\leq|v|_H$ and $|v|_J\leq|v|_H$.
Moreover $|D_av|_K\leq C\varepsilon^{-1}|v|_H$.
Every first-order term above is consequently bounded, when tested
against $v,w$, by $C\varepsilon^{-1}|v|_H|w|_H$; every mixed term
is bounded by $C\varepsilon^{-2}|v|_H|w|_H$.
The same holds for the derivatives of $J$.  Equivalently,
\begin{equation}\label{ana:graph-relative}
 \|H^{-1/2}H_aH^{-1/2}\|\leq C\varepsilon^{-1},
 \qquad
 \|H^{-1/2}H_{a\bar b}H^{-1/2}\|
 \leq C\varepsilon^{-2}.
\end{equation}

The connection matrix is $H^{-1}\partial H$.  In the order
$d\bar z_b\wedge dz_a$, its curvature coefficient is
$\partial_{\bar b}(H^{-1}H_a)$.  Conjugating by $H^{1/2}$, which
computes the operator norm in $H$, gives
\[
 H^{-1/2}H_{a\bar b}H^{-1/2}
 -(H^{-1/2}H_{\bar b}H^{-1/2})
  (H^{-1/2}H_aH^{-1/2}).
\]
Both terms have norm $O(\varepsilon^{-2})$ by
\eqref{ana:graph-relative}.  Passing to unitary tangent coordinates
costs only a fixed atlas constant.  Taking a maximum over the finitely
many stages gives $B$.  If $R_{a\bar b}$ are the Hermitian curvature
coefficients in such coordinates, then
\[
 \left|\sum_{a,b}\langle R_{a\bar b}\xi_a,\xi_b\rangle_H\right|
 \leq B\varepsilon^{-2}\left(\sum_a|\xi_a|_H\right)^2
 \leq NB\varepsilon^{-2}\sum_a|\xi_a|_H^2.
\]
No lower bound on a metric eigenvalue in a fixed Euclidean frame is
used in this calculation.
\end{proof}

\begin{proposition}\label{ana:homogeneous-derivative}
For each original homogeneous coordinate $Z_j$, the differentiated
row defines a bundle map
\[
 \partial_{Z_j}\Phi:E_0\longrightarrow L^{-1},
\]
and
\begin{equation}\label{ana:row-bound}
 |(\partial_{Z_j}\Phi)v|_{L^{-1}}
 \leq C\varepsilon^{-1}|v|_{H_{0,\varepsilon}}.
\end{equation}
The target has its original Fubini--Study line metric.
\end{proposition}

\begin{proof}
Put $D=D_0$.  If $D=0$ the derivative vanishes.  Otherwise, in the
chart $Z_k\neq0$ write $w_\ell=Z_\ell/Z_k$ and
$\Phi(Z)=Z_k^D\Phi_k(w)$.  The row $\partial_{Z_j}\Phi$ is
$Z_k^{D-1}$ times
\[
 P_{j,k}(w)=
 \begin{cases}
  \partial_{w_j}\Phi_k(w),&j\neq k,\\
  D\Phi_k(w)-\displaystyle\sum_{\ell\neq k}
       w_\ell\partial_{w_\ell}\Phi_k(w),&j=k.
 \end{cases}
\]
The coordinates are bounded on the smaller charts.
Proposition~\ref{ana:patched-derivatives}, together with
$\widehat G_0\leq H_0$ and $\varepsilon\leq1$, bounds each of these
rows by $C\varepsilon^{-1}$ in the indicated coefficient metric.
Its homogeneous degree is $D-1$, so its target as a map out of
$V_0\otimes L^{-D}$ is exactly $L^{-1}$.
The squared norm of the local target frame is $|Z_k|_L^{-2}$,
bounded on the chosen smaller chart.  This proves
\eqref{ana:row-bound}.  Equivalently, after tensoring by $L^{n+1}$,
the ratio of output and input line weights is this same fixed factor;
it is independent of $n$.
\end{proof}

\begin{proof}[Proof of Theorem~\ref{thm:D}]
The relative Gram estimates are proved in Lemma~\ref{ana:relative-gram}.
Propositions~\ref{ana:fiber-comparison} and~\ref{ana:patched-derivatives}
show that the chart metrics patch with the same bounds on actual fibers
on each smaller buffered chart $V_\alpha$.
The graph adjustment in Proposition~\ref{ana:graph-curvature} gives
contractivity, equivalence of metrics, and the asserted curvature bound
including the Nakano estimate. Proposition~\ref{ana:homogeneous-derivative}
proves the ambient homogeneous derivative estimate, with its indicated
line-bundle target. This proves all the assertions in the same metrics.
\end{proof}

\section{Microscopic division and a positive-twist estimate}
\label{ana:local-and-global}

The preceding metrics serve two purposes.  They normalize local
division uniformly, and their curvature can be dominated by a large
positive line twist.  We first keep the line power and the sampling
radius independent so that the order in which constants are chosen
remains visible.

\subsection{The Dolbeault solution estimate}

\begin{proposition}\label{ana:positive-gap}
There are $c_L>0$ and $C_0,C_1\geq0$, independent of the coefficient
stage, such that the following holds.  Let $h\geq0$ be an integer,
$0<\varepsilon\leq\varepsilon_0$, and
\[
 \Delta(h,\varepsilon)
 =c_Lh-C_0-C_1\varepsilon^{-2}>0.
\]
For every distributionally closed $L^2$ form $\zeta$ of type $(0,j)$,
$1\leq j\leq N$, with values in $E_q\otimes L^h$, there is a
unique solution perpendicular to $\ker\bar\partial$ in degree $j-1$,
and it satisfies
\begin{equation}\label{ana:solver}
 \bar\partial u=\zeta,
 \qquad
 \|u\|\leq\Delta(h,\varepsilon)^{-1/2}\|\zeta\|.
\end{equation}
All norms use $H_{q,\varepsilon}$, the original metric on $L^h$,
and the Fubini--Study form metric and volume.  At fixed parameters
the solution operator is linear and sends smooth closed forms to
smooth forms.
\end{proposition}

\begin{proof}
Put $F=E_q\otimes L^h$.  The canonical identification of $F$-valued
$(0,j)$-forms with $(F\otimes K_W^{-1})$-valued $(N,j)$-forms is
holomorphic and isometric when the canonical line has the induced
metric.  Locally one tensors a form with a nonvanishing holomorphic
volume form and its dual; their frame changes and metric factors
cancel.  Choose the order of the factors so that this identification
commutes with $\bar\partial$.  Its Hilbert adjoints then correspond
as well.

The Bochner--Kodaira--Nakano identity in top holomorphic degree gives
\[
 \|\bar\partial v\|^2+\|\bar\partial^*v\|^2
 \geq\int_W
 \langle[i\Theta(F\otimes K_W^{-1}),\Lambda]v,v\rangle\,dV.
\]
We use the sign convention in which the curvature eigenvalues of
the positive line $L$ are positive.  This identity and its $L^2$
consequences are recalled in \cite{Demailly,DemaillyL2}.
Let $c_L$ be a positive lower bound for the eigenvalues of $i\Theta(L)$.
On a top-holomorphic-degree form with antiholomorphic index set $J$,
the line contribution is
$h\sum_{a\in J}\lambda_a\geq hjc_L$.
The fixed curvature of $K_W^{-1}$ contributes a bounded operator,
which we allow to cost $C_0$.

For the coefficient curvature, fix an increasing $(j-1)$-index $S$
and denote by $v_{aS}$ the signed coefficient obtained by inserting
$a$ before $S$, with value zero when an index repeats.  The
top-degree curvature formula is the sum, over $S$, of the Nakano
forms of $(v_{aS})_{a=1}^N$.  By
Proposition~\ref{ana:graph-curvature} its absolute value is bounded
by
\[
 NB\varepsilon^{-2}\sum_{S,a}|v_{aS}|^2
 =NjB\varepsilon^{-2}|v|^2.
\]
Each increasing $j$-index occurs exactly $j$ times in the middle
sum.  We can therefore take $C_1=N^2B$, since $j\leq N$, and obtain
\begin{equation}\label{ana:basic-estimate}
 \|\bar\partial v\|^2+\|\bar\partial^*v\|^2
 \geq\Delta(h,\varepsilon)\|v\|^2.
\end{equation}

For completeness, on this compact manifold one may obtain the
solution directly from the Dolbeault Laplacian.  The smooth positive
metrics at fixed $(h,\varepsilon)$ make that operator elliptic and
self-adjoint.  Estimate \eqref{ana:basic-estimate} removes its kernel
in positive degree and bounds its inverse $G_j$ by $\Delta^{-1}$.
For closed data the Hodge identities give
\[
 u=\bar\partial^*G_j\zeta,
 \qquad \bar\partial u=\zeta.
\]
For smooth data these follow from commutation with the Laplacian;
$\bar\partial G_j\zeta$ is harmonic and hence zero in the next
degree, or is zero by degree at $j=N$.  The same identities extend
to $L^2$ closed data by the spectral calculus, or by smooth heat
regularization, which commutes with $\bar\partial$.  Moreover
\[
 \|u\|^2
 =\langle\zeta,G_j\zeta\rangle
 \leq\Delta^{-1}\|\zeta\|^2.
\]
The image of $\bar\partial^*$ is perpendicular to the closed forms.
Any two solutions differ by a closed form, proving minimality and
uniqueness.  These also imply linearity.  Elliptic regularity gives
smoothness for smooth inputs.
\end{proof}

\subsection{Division on shrinking balls with the original line weight}

Write $B(p,r)$ for a geodesic ball in $W$.  For a section of
$E_q\otimes L^h$ over $U$, let $\|\cdot\|_{q,h;U}$ denote its
$L^2$ norm in $H_{q,\varepsilon}\otimes L^h$.  At $q=-1$ this is
the original scalar line norm.

\begin{proposition}\label{ana:microscopic-division}
For every $s>0$ there are $\Lambda(s)>s$ and
$\varepsilon_0(s)>0$ with the following property.  Fix $K\geq0$.
Suppose $h\geq0$ is an integer,
$0<\varepsilon\leq\varepsilon_0(s)$, and $h\varepsilon^2\leq K$.
Let $b$ be a square-integrable holomorphic section of
$E_{q-1}\otimes L^h$ on $B(p,\Lambda(s)\varepsilon)$.
For $q=0$ assume $b$ is a section of $\cI\otimes L^h$ there;
for $1\leq q\leq N$ assume $d_{q-1}b=0$ there.
There is a holomorphic section $a$ on $B(p,s\varepsilon)$ such that
\begin{equation}\label{ana:micro-bound}
 d_qa=b,
 \qquad
 \|a\|_{q,h;B(p,s\varepsilon)}
 \leq C_s e^{C\Lambda(s)^2K}
 \|b\|_{q-1,h;B(p,\Lambda(s)\varepsilon)}.
\end{equation}
The choice can be made linear in $b$ for each fixed set of parameters.
The constants are uniform in the center, the radius, and the line power.
\end{proposition}

\begin{proof}
If $E_q=0$, the image condition and the exactness at the preceding
stage give $b=0$ and the lift is zero. We may thus assume $E_q\ne0$.
Choose $\ell>0$ so that every geodesic ball $B(p,2\ell)$ is
contained in some smaller chart $V_\alpha$; this follows from a
Lebesgue number for the finite smaller-chart cover. The compact
buffers in $U_\alpha$ and the finitely many smooth coordinate maps
give, after decreasing $\ell$ if necessary, a common constant
$H\geq1$ such that for $0<r\leq\ell$ in such a chart,
\[
 B_{\rm coord}(p,r/H)\subset B(p,r)\subset B_{\rm coord}(p,Hr),
\]
where $B_{\rm coord}$ denotes a Euclidean coordinate ball. These
inclusions follow from the uniform upper and lower bounds for the
Fubini--Study metric in the buffered coordinates; one may take the
radii sufficiently small that the relevant straight segments and
minimizing geodesics remain in $U_\alpha$.

Proposition~\ref{alg:compact} and Theorem~\ref{alg:hilbert} supply
compact families of normalized homogeneous matrices with fixed image
dimensions in every degree, after the common target shift is removed.
For output factor $Hs$, let $\lambda_q(s)>Hs$ be the input factor in
Theorem~\ref{div:uniform} for the family at stage $q$, and take
$\Lambda(s)=H\max_{0\leq q\leq N}\lambda_q(s)>s$, assigning any larger
factor than $Hs$ at a zero stage. Decrease $\varepsilon_0(s)$ so
that $\Lambda(s)\varepsilon<\ell$. For each $p$ choose the smaller
chart containing $B(p,2\ell)$. The inclusions now give
\[
 B(p,s\varepsilon)\subset B_{\rm coord}(p,Hs\varepsilon)
 \subset B_{\rm coord}(p,\lambda_q(s)\varepsilon)
 \subset B(p,\Lambda(s)\varepsilon)\subset V_\alpha.
\]
In particular all the balls on which division is performed lie in the
domain of the uniform comparison in Theorem~\ref{thm:D}. In this
chart freeze the Gram metrics at $p$ and put
$R_j=G_{j,\alpha,\varepsilon}(p)^{1/2}$, $R_{-1}=1$.
In the rescaled variable $\zeta=(x-p)/\varepsilon$, the matrix is
\[
 A_q(\zeta)=R_{q-1}d_{q,\alpha}(p+\varepsilon\zeta)R_q^{-1}.
\]
This matrix belongs to the family just used to choose $\lambda_q(s)$.

The image condition holds on the entire input ball. At $q=0$ it is
the assumed ideal-sheaf membership. At $1\leq q\leq N$ it follows
from exactness of the actual presentation at $E_{q-1}$, since $b$
is a holomorphic cycle. Analytification preserves this exactness.
We use membership in the analytic image sheaf, not merely membership
in the image of an evaluated matrix at each point. Restrict the datum
to the coordinate input ball, apply the normalized division theorem,
and restrict its output to the desired geodesic ball.

In a holomorphic frame of the common factor $L^h$, the normalized
datum is $\widetilde b(\zeta)=R_{q-1}b(p+\varepsilon\zeta)$.
Undoing the normalization of the primitive gives a local section of
$E_q\otimes L^h$.  The change of volume contributes the same
factor $\varepsilon^N$ to both unsquared norms.  The chart volume
densities are uniformly bounded above and below.
Lemma~\ref{ana:relative-gram} compares the frozen and moving Gram
norms on these fixed multiples of $\varepsilon$.
Propositions~\ref{ana:fiber-comparison} and
\ref{ana:graph-curvature} compare them with $H_q$ and $H_{q-1}$.
All resulting constants depend on the selected radii, but not on
$p$, $\varepsilon$, or $h$.

It remains to restore the common line weight.  Write
$|e(x)|_L^2=e^{-\varphi(x)}$ in a local holomorphic frame.  Replace
that frame, for this local calculation, by
\[
 e'(x)=
 \exp\!\left(\frac12\varphi(p)
        +\sum_a\partial_a\varphi(p)(x_a-p_a)\right)e(x).
\]
Its potential has zero constant and linear Taylor terms at $p$;
hence its absolute value is at most $C|x-p|^2$ on the relevant
buffered chart.  Use $(e')^h$ as the frame of the \emph{common}
factor $L^h$, while keeping the coefficient-bundle frames fixed.
Each polynomial matrix $d_{q,\alpha}$ is unchanged, because the
same extra line frame occurs in its source and target.
The squared line weight and its reciprocal are bounded on the
input ball by $e^{C\Lambda(s)^2h\varepsilon^2}$.
Comparing the two norms proves \eqref{ana:micro-bound}, after a
change of the fixed constant $C$.

For a fixed matrix, Theorem~\ref{div:uniform} permits the
minimum-norm primitive on its coordinate output ball.  That choice
is linear on admissible data.  Normalization, line-frame changes,
and restrictions preserve linearity.  No continuity of this choice
as $p$ or $\varepsilon$ varies is required.
\end{proof}

\section{Holomorphic overlap data and the global lift}
\label{ana:global-section}

We use local division only on holomorphic sections.  A shrinking
cover converts these holomorphic overlap data into smooth Dolbeault
forms.  The loss from each differentiated cutoff is then recovered
by one use of the positive-twist solution estimate.

\begin{thmE}[Bounded projective lifting]\label{geo:lifting}
Let $\cI\subset\cO_W$ be a nonzero proper coherent ideal sheaf on
$W=\Pj^N$, $N\ge1$.  Let $D=D_0$ be chosen as in Theorem~\ref{thm:B}, with complete
space $V=H^0(W,\cI\otimes L^D)$ and evaluation map
\[
 \Phi:E_0=V\otimes L^{-D}\longrightarrow\cO_W,
 \qquad\im\Phi=\cI.
\]
Use the metrics of Theorem~\ref{thm:D}, in particular
\eqref{eq:geo-contract}--\eqref{eq:geo-curvature}.
There are fixed constants $R_*>0$, $n_0$, and $C'$ such that for
$n\ge n_0$, $\varepsilon=R_*/\sqrt n$, and $h\in\{n,n+1\}$,
every $F\in H^0(W,\cI\otimes L^h)$ has a global holomorphic lift
$a\in H^0(W,E_0\otimes L^h)$ satisfying
\begin{equation}\label{eq:geo-lift-bound}
 \Phi a=F,\qquad
 \|a\|_{H_{0,\varepsilon}\otimes L^h,L^2}
       \le C'\|F\|_{L^h,L^2}.
\end{equation}
The right-hand norm is the original scalar Fubini--Study norm.
The lift can be chosen linearly in $F$.  All constants and the
presentation are independent of $n,h,F$.  The same coefficient metric
serves the two adjacent powers. The unit ideal sheaf has the same lifting
conclusion with the row $1$ and the constant metric.
\end{thmE}

The derivatives in \eqref{eq:geo-differentiate} are derivatives of the
homogeneous polynomial entries of $\Phi$ in the fixed ambient
coordinates.  Their degree is $D-1$, so they define bundle maps
$E_0\to L^{-1}$, as indicated.  For $D=0$ they vanish.

\paragraph{Historical relation of the descent.}
The finite correction through a complex is a version of the
Koszul--Dolbeault procedure in \cite{KelleherTaylor,HormanderGenerators}
and of the projective construction in \cite[Section~6]{AnderssonWulcan}.
In the latter the residue-current identities supply data for successive
Dolbeault equations. Here the cochains are holomorphic overlap lifts,
obtained with uniform estimates from Theorem C before cutoffs are
introduced. We prove the precise norm bounds and closedness at every
stage; no vanishing theorem alone would imply the quantitative conclusion.

\subsection{A cover with buffers for every face}

Choose a maximal $\varepsilon$-separated set of centers
$\{p_i\}$ and fix an ordering.  The balls $B(p_i,\varepsilon)$
cover $W$.  For any fixed $T>0$, there is a bound $M(T)$,
independent of sufficiently small $\varepsilon$, for the number of
centers within $T\varepsilon$ of an arbitrary point.  Indeed the
balls of radius $\varepsilon/2$ about these centers are disjoint
and lie in a ball of radius $(T+1/2)\varepsilon$.
Uniform upper and lower small-ball volume estimates on the compact
manifold give, for example, $M(T)\leq C(1+2T)^{2N}$.

Fix $\kappa>1$.  Smooth bumps equal to one on
$B(p_i,\varepsilon)$ and supported strictly inside
$B(p_i,\kappa\varepsilon)$ can be obtained from a fixed function of
$\dist(p_i,x)^2/\varepsilon^2$ below the injectivity radius.
Their first derivatives are bounded by $C/\varepsilon$.
Dividing them by their sum gives a partition of unity $\{\chi_i\}$
with
\begin{equation}\label{ana:shrinking-cutoffs}
 \sum_i\chi_i=1,\qquad
 \supp\chi_i\Subset B(p_i,\kappa\varepsilon),\qquad
 |\bar\partial\chi_i|\leq C\varepsilon^{-1}.
\end{equation}
At every point only a uniformly bounded number of terms contributes.
These cutoffs vary with the scale; the atlas partition in
\eqref{ana:patched-metric} was fixed earlier.

For $0\leq q\leq N$, call an increasing tuple
$I=(i_0,\ldots,i_q)$ of degree $q$ \emph{relevant} if the closed
supports of its cutoffs have a common point. The cover may admit
higher-order intersections, but no tuples beyond degree $N$ are
needed or assigned balls or coefficient cochains.  Every nonempty face
of a relevant tuple is relevant.  All its centers lie within
$2\kappa\varepsilon$ of one another.  Anchor $I$ at its smallest
index, and use that same anchor for every permutation of $I$.

Choose $A_N>2\kappa+2$ and then, for $q=N,\ldots,1$, choose
\begin{equation}\label{ana:buffer-recursion}
 A_{q-1}>\Lambda(A_q)+2\kappa+2.
\end{equation}
For a relevant tuple $I$ of degree $0\leq q\leq N$ let
\[
 V_I=B(p_{i_0},A_q\varepsilon),
 \qquad
 U_I=B(p_{i_0},\Lambda(A_q)\varepsilon).
\]

\begin{lemma}\label{ana:buffers}
If $J$ is a nonempty proper face of a relevant tuple $I$ of degree
$1\leq q\leq N$, then $U_I\subset V_J$.  A common $\varepsilon_{\rm buf}>0$
makes all these input and output balls admissible in
Proposition~\ref{ana:microscopic-division}, uniformly in the centers.
\end{lemma}

\begin{proof}
First suppose that $J$ is obtained from $I$ by deleting one index.
The two anchors have distance at most $2\kappa\varepsilon$.
For $x\in U_I$, therefore,
\[
 \dist(x,p_J)
 <\bigl(\Lambda(A_q)+2\kappa\bigr)\varepsilon
 <(A_{q-1}-2)\varepsilon.
\]
This includes the face obtained by deleting the smallest index,
whose anchor changes.  For a nonempty proper face of higher codimension, choose a chain of
successive single-index deletions from $I$ to that face. Applying the
codimension-one inclusion at each step and using $V_K\subset U_K$ for
each intermediate tuple $K$ gives the stated inclusion.  There are only finitely many
fixed radii $A_q$ and $\Lambda(A_q)$.  Decreasing the admissible
scale places every ball in a suitable buffered chart and makes
all the preceding coordinate, volume, and cutoff estimates hold.
\end{proof}

\subsection{Fixing the scale}

At this point the presentation, metric estimates, local division
radii, and buffers have all been fixed.  In particular the constants
$c_L,C_0,C_1$ of Proposition~\ref{ana:positive-gap} do not depend
on any norm of a division operator.  Choose a fixed $R_*>0$ with
\begin{equation}\label{ana:R-choice}
 C_1/R_*^2<c_L/2.
\end{equation}
Increase an integer $n_0$ so that
\[
 n_0\geq1,\qquad
 R_*/\sqrt{n_0}\leq\varepsilon_{\rm buf},\qquad
 C_0\leq c_Ln_0/4.
\]
We may also require $n_0\geq D_0$.  For $n\geq n_0$ put
\begin{equation}\label{ana:scale}
 \varepsilon=R_*/\sqrt n.
\end{equation}
For both $h=n$ and $h=n+1$,
\begin{equation}\label{ana:simultaneous-scale}
 h\varepsilon^2\leq2R_*^2,
 \qquad
 \Delta(h,\varepsilon)\geq c_Ln/4.
\end{equation}
Local division is consequently uniform at all prescribed radii,
and the minimal global solution satisfies
\begin{equation}\label{ana:scaled-solver}
 \|u\|\leq\frac{2}{\sqrt{c_Ln}}\|\zeta\|
 =b\varepsilon\|\zeta\|,
 \qquad b=\frac{2}{R_*\sqrt{c_L}}.
\end{equation}
The constants for local division may now contain factors such as
$\exp(2C\Lambda(A_q)^2R_*^2)$.  These are fixed constants.
They do not enter the earlier curvature constant $C_1$, and the
finite correction below imposes no smallness condition on them.

\subsection{Holomorphic cochains with uniform total norm}

Fix $n\geq n_0$, one of the two powers $h\in\{n,n+1\}$, and
$F\in H^0(W,\cI\otimes L^h)$.  Write $\|F\|$ for its original
scalar $L^2$ norm.  For a cochain of degree $q-1$, define its
alternating difference on an increasing relevant tuple $I$ of degree
$1\leq q\leq N$ by
\[
 (\delta c)_I
 =\sum_{\ell=0}^q(-1)^\ell
 c_{I\setminus i_\ell}|_{U_I}.
\]
Lemma~\ref{ana:buffers} makes the restrictions meaningful.
These are differences of sections of the same actual bundle;
columns in two chart frames must first be identified as such sections.

\begin{proposition}\label{ana:holomorphic-cochains}
For every increasing relevant tuple $I$ of degree $0\leq q\leq N$, there is a
holomorphic section $c_{q,I}$ of $E_q\otimes L^h$ on $V_I$, linear
in $F$, with
\begin{equation}\label{ana:cochain-identities}
 \Phi c_{0,i}=F|_{V_i},
 \qquad
 d_qc_{q,I}=(\delta c_{q-1})_I|_{V_I}\quad(1\leq q\leq N).
\end{equation}
Moreover
\begin{equation}\label{ana:cochain-mass}
 S_q:=\sum_{I\ \mathrm{increasing,\ relevant}}
       \|c_{q,I}\|_{q,h;V_I}^2
 \leq C_q\|F\|^2,
\end{equation}
with constants independent of $n$, $h$, and $F$.
\end{proposition}

\begin{proof}
At degree zero divide $F$ on $U_i$ using
Proposition~\ref{ana:microscopic-division}.  At degree one,
$c_{0,j}-c_{0,i}$ is a holomorphic section on $U_{ij}$ whose image
under $\Phi$ is $F-F=0$ on that entire ball.  It belongs to the
image sheaf of $d_1$, so division gives $c_{1,ij}$ on $V_{ij}$.

Suppose the construction has reached degree $q-1$, with $2\leq q\leq N$.
All face identities hold on the whole new input $U_I$, by the
buffer inclusions, including those for faces of faces.  Hence
\[
 d_{q-1}(\delta c_{q-1})_I
 =(\delta^2c_{q-2})_I=0\quad\hbox{on }U_I.
\]
Indeed the two orders of deleting positions $r<s$ have signs
$(-1)^{r+s-1}$ and $(-1)^{r+s}$ and restrict the same section to
the same domain. Exactness at $E_{q-1}$ supplies analytic image-sheaf
membership of this holomorphic datum throughout $U_I$. This is within
the exactness range because $q\leq N$.  Local division gives
$c_{q,I}$ on $V_I$.  Fix all charts and the linear local division
operators before varying $F$; induction proves linear dependence.
For permutations of an increasing tuple use the corresponding sign
and the original anchored ball, and put the cochain equal to zero
on tuples with repeated indices.

Let $L_q$ be a common local-division norm at output factor $A_q$
for the fixed bound $K=2R_*^2$ on $h\varepsilon^2$.  Every point belongs to at most
$M(\Lambda(A_0))$ initial input balls.  Thus
\[
 S_0\leq L_0^2M(\Lambda(A_0))\|F\|^2.
\]
At degree $1\leq q\leq N$, the squared norm of the alternating sum is
bounded by $q+1$ times the sum of the squared face norms on $U_I$.
Each is at most the full norm on $V_J$.  A fixed nonempty face $J$
has at most $M(2\kappa)$ relevant increasing cofaces: the new
index must have its center within $2\kappa\varepsilon$ of the
anchor of $J$.  Consequently
\[
 S_q\leq L_q^2(q+1)M(2\kappa)S_{q-1}.
\]
Iterating at most $N$ times proves \eqref{ana:cochain-mass}.
The number of balls, which grows as the scale decreases, never
enters this bound.
\end{proof}

\subsection{The cochain-to-form identities}

Set $\alpha_i=\bar\partial\chi_i$.  Define, with ordered tuples,
\begin{equation}\label{ana:eta-definition}
 \eta_q=
 \sum_{i_0,\ldots,i_q}
 \chi_{i_0}\alpha_{i_1}\wedge\cdots\wedge\alpha_{i_q}
 c_{q,i_0\ldots i_q},\qquad 0\leq q\leq N.
\end{equation}
For $q=0$ this means $\eta_0=\sum_i\chi_ic_{0,i}$.
An irrelevant tuple contributes zero.  The common cutoff support
is compactly contained in $V_I$, since $A_q>2\kappa+2$.
Every summand therefore extends by zero to a smooth global form.

\begin{proposition}\label{ana:eta-identities}
These forms have values in $E_q\otimes L^h$ and satisfy
\begin{equation}\label{ana:eta-complex}
 \Phi\eta_0=F,
 \qquad
 d_q\eta_q=\bar\partial\eta_{q-1}\quad(1\leq q\leq N),
 \qquad
 \bar\partial\eta_N=0.
\end{equation}
There are fixed constants $P_q$ such that
\begin{equation}\label{ana:eta-estimate}
 \|\eta_q\|_{H_{q,\varepsilon}\otimes L^h}
 \leq P_q\varepsilon^{-q}\|F\|,
 \qquad 0\leq q\leq N.
\end{equation}
All forms depend linearly on $F$.
\end{proposition}

\begin{proof}
The first identity follows from $\sum_i\chi_i=1$ and
\eqref{ana:cochain-identities}.  For the second, apply $d_q$ to
\eqref{ana:eta-definition} and expand the alternating difference.
The face deleting $i_0$ gives
\[
 \sum_{i_1,\ldots,i_q}
 \alpha_{i_1}\wedge\cdots\wedge\alpha_{i_q}
 c_{q-1,i_1\ldots i_q},
\]
because the sum over $i_0$ of $\chi_{i_0}$ is one.  This is
exactly $\bar\partial\eta_{q-1}$: the local coefficient sections
are holomorphic and $\bar\partial\alpha_i=0$.
For a face deleting $i_\ell$ with $\ell\geq1$, its coefficient is
independent of that index.  Summing the factor $\alpha_{i_\ell}$
in its unchanged wedge position gives zero, since
$\sum_i\alpha_i=0$.  These calculations are pointwise legitimate:
every tuple whose cutoff factors contribute at a point is relevant,
and every coefficient used there is defined on a neighborhood of
that point.  Repeated-index terms obey the alternating convention.
There is no factorial prefactor in \eqref{ana:eta-definition} or
in these identities.  Finally a $(0,N+1)$-form on $W$ vanishes,
which proves the last identity in \eqref{ana:eta-complex} without any
injectivity assumption on $d_N$.

For the norm estimate, at a point at most $M(\kappa)^{q+1}$
ordered tuples can contribute.  The norm of an individual summand
is at most $C_q\varepsilon^{-q}|c_{q,I}|$ on its support, by
\eqref{ana:shrinking-cutoffs} and the exterior-product inequality.
Apply the squared Cauchy--Schwarz inequality to the finite sum and
integrate.  For each increasing relevant tuple there are at most
$(q+1)!$ ordered versions, all using the same anchored section and
ball.  Thus
\[
 \|\eta_q\|^2
 \leq C_q M(\kappa)^{q+1}(q+1)!
        \varepsilon^{-2q} S_q.
\]
Proposition~\ref{ana:holomorphic-cochains} proves
\eqref{ana:eta-estimate}, also at $q=0$, where no differentiated
cutoff appears.  Linearity follows from the construction with
fixed cutoffs.
\end{proof}

\subsection{Finite descending correction}

At a fixed $(n,h)$ all the following spaces carry the same coefficient
metrics and the same line twist.  Begin with $\zeta_N=\eta_N$.
For $q=N,\ldots,1$, solve in the coefficient bundle $E_q\otimes L^h$
\begin{equation}\label{ana:descent}
 \bar\partial u_q=\zeta_q,
 \qquad
 \zeta_{q-1}=\eta_{q-1}-d_qu_q.
\end{equation}
The form $u_q$ has type $(0,q-1)$, so the two terms defining
$\zeta_{q-1}$ have the same bundle and form degree.

\begin{lemma}\label{ana:descent-closed}
Every input $\zeta_q$ in \eqref{ana:descent} is closed.
The procedure gives a holomorphic section $a=\zeta_0$ with
$\Phi a=F$.
\end{lemma}

\begin{proof}
The top form is closed by \eqref{ana:eta-complex}.  If $\zeta_q$
is closed, Proposition~\ref{ana:positive-gap} and
\eqref{ana:simultaneous-scale} give its smooth minimal solution.
The presentation maps are holomorphic and commute with
$\bar\partial$, so
\[
 \bar\partial\zeta_{q-1}
 =d_q\eta_q-d_q\bar\partial u_q
 =d_q(\eta_q-\zeta_q).
\]
For $q=N$ the last difference is zero.  For $q<N$, the preceding
step says $\eta_q-\zeta_q=d_{q+1}u_{q+1}$, and the expression
vanishes because $d_qd_{q+1}=0$.  This proves closedness at every
step, and hence justifies each subsequent application of the solver.
At degree zero, closedness means that $a$ is holomorphic.  Finally
\[
 \Phi a=\Phi\eta_0-\Phi d_1u_1=F.
\]
If a coefficient stage is zero, the corresponding data and solutions
are zero; the same induction applies.
\end{proof}

\begin{lemma}\label{ana:descent-norm}
Let $B_N=P_N$ and define the finite recursion
\[
 B_{q-1}=P_{q-1}+bB_q,
 \qquad q=N,\ldots,1,
\]
where $b$ is from \eqref{ana:scaled-solver}.  Then
\[
 \|\zeta_q\|\leq B_q\varepsilon^{-q}\|F\|,
 \qquad
 \|u_q\|\leq bB_q\varepsilon^{-(q-1)}\|F\|.
\]
In particular,
\begin{equation}\label{ana:final-lift-bound}
 \|a\|_{H_{0,\varepsilon}\otimes L^h}
 \leq \left(\sum_{q=0}^N b^qP_q\right)\|F\|_{L^h}.
\end{equation}
\end{lemma}

\begin{proof}
The top estimate is \eqref{ana:eta-estimate}.  At a closed input,
\eqref{ana:scaled-solver} returns one factor $\varepsilon$ and
gives the bound for $u_q$.  Contractivity of $d_q$ in the graph
metrics, preserved on line-twisted forms, then gives
\begin{align*}
 \|\zeta_{q-1}\|
 &\leq\|\eta_{q-1}\|+\|d_qu_q\|\\
 &\leq(P_{q-1}+bB_q)\varepsilon^{-(q-1)}\|F\|.
\end{align*}
This proves the induction.  Expanding its finite recursion gives
$B_0=\sum_{q=0}^N b^qP_q$.  In particular no contraction condition
on $b$ or on a local-division constant is needed.  Closedness was
proved by the complex identities, so no estimate for a derivative
of a corrected form is being used in this norm argument.
\end{proof}

\begin{proposition}\label{ana:global-lift}
For all $n\geq n_0$ and $h\in\{n,n+1\}$ there is a linear map
\[
 T_{n,h}:H^0(W,\cI\otimes L^h)
 \longrightarrow H^0(W,E_0\otimes L^h)
\]
such that $\Phi T_{n,h}=\id$ and
\[
 \|T_{n,h}F\|_{H_{0,R_*/\sqrt n}\otimes L^h}
 \leq C\|F\|_{L^h}.
\]
The datum on the right has its original scalar Fubini--Study norm.
The same coefficient metric is used at the two adjacent line powers.
\end{proposition}

\begin{proof}
For fixed $(n,h)$ fix the cover, partition, frames, and local
minimum-norm division operators.  Proposition~\ref{ana:holomorphic-cochains}
and formula \eqref{ana:eta-definition} depend linearly on $F$.
Each minimal Dolbeault solution is linear on the closed subspace of
admissible inputs by Proposition~\ref{ana:positive-gap}.
Lemma~\ref{ana:descent-closed} ensures that every corrected input
belongs to that subspace.  The finite recursion consequently gives
a linear $a(F)$.  Set $T_{n,h}F=a(F)$.  Its right-inverse identity
and holomorphy follow from Lemma~\ref{ana:descent-closed}, and its
norm bound is \eqref{ana:final-lift-bound}.  All constants in that
bound were fixed before $n$ varies, and
\eqref{ana:simultaneous-scale} holds at both powers.
The two linear operators may be chosen separately; no identification
of their source spaces is asserted.
\end{proof}

\subsection{Completion of the lifting theorem}

\begin{proof}[Proof of Theorem~\ref{geo:lifting}]
For a proper nonzero ideal sheaf in positive projective dimension,
the complete presentation of Proposition~\ref{alg:complete} supplies
the fixed row $\Phi$ and degree $D=D_0$.  The graph metrics
\eqref{ana:graph-definition} are smooth and positive at every
admissible radius, and make $\Phi$ contractive into the original
scalar norm.  Proposition~\ref{ana:homogeneous-derivative} proves
the homogeneous derivative estimate with target $L^{-1}$ in that
same metric.  Proposition~\ref{ana:graph-curvature} proves its
$O(\varepsilon^{-2})$ curvature bound, measured in its own fiber
norm.  Finally Proposition~\ref{ana:global-lift} proves the bounded
holomorphic lifting assertion at $\varepsilon=R_*/\sqrt n$, for
both $h=n$ and $h=n+1$, with the original norm on the datum.
Taking the maximum of the finitely many constants gives the single
constant permitted in the statement.

For the unit ideal take $D=0$, $V=\C$, $\Phi=1$, and the constant
coefficient metric.  The identity lift has norm one and the row
derivative and curvature vanish.  For completeness, on $\Pj^0$ every nonzero ideal
sheaf is the unit ideal, so its lift is also the identity.
The zero ideal requires no lifting: its discarded coordinate
components vanish.  This proves every asserted clause.
\end{proof}

\begin{remark}\label{ana:constant-order}
The choices have a definite order.  The complete presentation fixes
the polynomial degrees.  Those degrees and the fixed atlas determine
the Gram estimates and the curvature constants.  Compact normalized
division determines an input radius for each output radius.  The
finite buffer recursion fixes the shrinking-cover geometry.
Only then is $R_*$ chosen to dominate the curvature loss.
The line-weight comparison changes the local-division constants
after this choice, but it does not change the curvature constants.
The finite descent accommodates those constants without any
smallness assumption.  This order accounts for all uniformities
in the lifting theorem.
\end{remark}

% END sections/analysis.tex


% BEGIN sections/conclusion.tex
\section{Step 3: Return to Drury--Arveson space}
\label{sec:completion}

The preceding construction proves Theorem~\ref{geo:lifting}. We now
return to the original scalar norm, obtain the operator estimate, and
transfer it to a vector module before the Hilbert function is used.

\subsection{The scalar estimate in the original Hilbert norm}
\begin{proposition}\label{geo:transfer}
Theorem~\ref{geo:lifting} implies the scalar degree estimate
\eqref{eq:main-degree} for every homogeneous ideal in every finite
number of variables.
\end{proposition}
\begin{proof}
The exceptional cases and the finite change under saturation have
already been treated.  Work in a degree where
$I_n=I_n^{\mathrm{sat}}$ and $I_{n+1}=I_{n+1}^{\mathrm{sat}}$.
Let $F\in I_{n+1}$ and use Theorem~\ref{geo:lifting} with $h=n+1$.
The entries of its global holomorphic lift $a$ are homogeneous
polynomials of degree $n+1-D$.  Differentiating the equality $F=\Phi a$
gives
\[
 \partial_{z_j}F=(\partial_{z_j}\Phi)a+\Phi\partial_{z_j}a.
\]
The second term belongs to $I_n^{\mathrm{sat}}=I_n$; if the
coefficient degree is zero it is simply zero.  The common orthogonal
projection in Lemma~\ref{geo:norm} therefore gives
\[
 P_{Q_n}\partial_{z_j}F
   =P_{Q_n}\bigl((\partial_{z_j}\Phi)a\bigr).
\]
Apply \eqref{eq:geo-differentiate} pointwise and
\eqref{eq:geo-lift-bound} after integration.  Since the projection is
contractive in the scalar norm,
\begin{equation}\label{eq:geo-projected-derivative}
 \|P_{Q_n}\partial_{z_j}F\|_{\FS,n}
       \le C''\varepsilon^{-1}\|F\|_{\FS,n+1}.
\end{equation}
On degree $n+1$, the Drury--Arveson adjoint is
$Z_j^*=(n+1)^{-1}\partial_{z_j}$, and
$B_{j,n}^*F=P_{Q_n}Z_j^*F$.  Consequently
\begin{align}
 \|B_{j,n}^*F\|_{\DA,n}
 &\le \frac{C''}{\varepsilon(n+1)}
       \sqrt{\frac{c_n}{c_{n+1}}}\|F\|_{\DA,n+1}\notag\\
 &=\frac{C''}{\varepsilon\sqrt{(n+1)(n+m)}}
       \|F\|_{\DA,n+1}.\label{eq:geo-exact-transfer}
\end{align}
At $\varepsilon=R_*/\sqrt n$, summing the squared operator bounds gives
\[
 A_n\le
 \frac{m(C'')^2n}{R_*^2(n+1)(n+m)}I_{Q_n}
 \le\frac{C_I}{n+1}I_{Q_n}.
\]
The finitely many omitted degrees are absorbed into $C_I$.
\end{proof}

No derivative of the metric or of the chosen lift occurs in the last
calculation.  The derivative of the lift has been removed by ideal
membership.  Curvature and derivatives of the metric are needed earlier,
to construct a holomorphic lift with a compatible norm bound.  This
distinction separates the analytic construction from its operator
consequence.


In particular, Proposition~\ref{geo:transfer}, with the finite-degree
adjustment of Lemma~\ref{geo:saturation}, proves for every homogeneous
scalar ideal, in any finite number of variables, the estimate
\begin{equation}\label{eq:complete-scalar}
 A_n^I\le\frac{C_I}{n+1}I_{Q_n}.
\end{equation}

\subsection{An auxiliary-variable realization}
Let $M\subset R\otimes\C^r$, where $R=\C[z_1,\ldots,z_m]$, and
let $T=(R\otimes\C^r)/M$.  Adjoin $r$ variables and define
\[
 R'=R[y_1,\ldots,y_r],\qquad
 \lambda(f)=\sum_{b=1}^r f_b(z)y_b\quad(f\in R\otimes\C^r).
\]
Choose homogeneous generators $f^{(1)},\ldots,f^{(s)}$ of $M$ and set
\begin{equation}\label{eq:complete-encoding}
 J_M=(\lambda(f^{(1)}),\ldots,\lambda(f^{(s)}))+(y_1,\ldots,y_r)^2.
\end{equation}
The auxiliary-variable realization and its degreewise normalization
occur in Shalit's stable-division argument
\cite[Section~5, equation~(5.1)]{ShalitDivision}.  The square-zero
relations in \eqref{eq:complete-encoding} let us record the complete
positive defect exactly.

\begin{lemma}\label{complete:ring}
The ideal $J_M$ is homogeneous separately in the $z$ and $y$ variables.
Its component of $y$-degree one is exactly $\lambda(M)$.  There is an
isomorphism of graded algebras
\[
 R'/J_M\simeq R\ltimes T(-1),
 \qquad (a,t)(a',t')=(aa',at'+a't).
\]
In particular $\sqrt{J_M}=(y_1,\ldots,y_r)$ and
$\dim R'/J_M=m$.
\end{lemma}
\begin{proof}
Each generator of \eqref{eq:complete-encoding} has definite degree in
both groups of variables.  In $y$-degree zero no relation occurs; in
$y$-degree one only coefficients in $R$ can multiply the first set of
generators.  Their image is exactly $\lambda(M)$ by $R$-linearity.
All higher $y$-degrees vanish.  Thus a class has the displayed pair
description, and multiplication has the claimed form.  The shift is
$T(-1)_\ell=T_{\ell-1}$, because each auxiliary variable adds one to
total degree.  Finally
$(y)^2\subset J_M\subset(y)$ and $(y)$ is prime, with quotient $R$.
\end{proof}

Let $Q^{J_M}\subset H_{m+r}^2$ be the scalar quotient and let $E_n$
be its sector of $z$-degree $n$ and $y$-degree one.  Its total degree is
$n+1$.  Write $Q_n^M$ for the original vector quotient degree.

\begin{lemma}\label{complete:unitary}
The map
\begin{equation}\label{eq:complete-unitary}
 U_n:Q_n^M\longrightarrow E_n,\qquad
 U_nq=\sqrt{n+1}\sum_{b=1}^r q_b(z)y_b
\end{equation}
is unitary.  In total degree $\ell\ge1$,
\[
 Q^{J_M}_\ell=H_{m,\ell}^2\oplus E_{\ell-1}
\]
orthogonally; in degree zero it is $\C$.
\end{lemma}
\begin{proof}
For $|\alpha|=n$,
\[
 \|z^\alpha e_b\|_{H_m^2\otimes\C^r}^2=\frac{\alpha!}{n!},
 \qquad
 \|z^\alpha y_b\|_{H_{m+r}^2}^2=\frac{\alpha!}{(n+1)!}.
\]
Distinct displayed monomials are orthogonal.  Therefore the same
formula defines a unitary $\widetilde U_n$ between the full vector
degree space and the full scalar sector linear in $y$.  It maps $M_n$
onto the corresponding sector of $J_M$ by Lemma~\ref{complete:ring},
and hence maps their orthogonal complements onto one another.
The other surviving sector is the pure $z$ sector, whose monomial norms
are unchanged.  All higher auxiliary degrees lie in the ideal.
\end{proof}

Let $A^{J_M}_\ell$ be the scalar defect square, summed over all
$m+r$ coordinates.  The following equality supplies the quantitative
transfer.

\begin{proposition}\label{complete:defect}
The operator $A^{J_M}_\ell$ preserves the two surviving auxiliary-degree
sectors.  On the linear sector,
\begin{equation}\label{eq:complete-defect}
 U_n^*A^{J_M}_{n+1}U_n
 =\frac{n+1}{n+2}A_n^M+\frac{r+1}{n+2}I_{Q_n^M}.
\end{equation}
On the pure $z$ sector in total degree $\ell$,
\[
 0\le A^{J_M}_\ell|_{H_{m,\ell}^2}
       \le\frac r{\ell+1}I.
\]
\end{proposition}
\begin{proof}
Let $q\in Q_n^M$.  For each original coordinate,
\[
 z_jU_nq
 =\sqrt{\frac{n+1}{n+2}}\,\widetilde U_{n+1}(z_jq).
\]
The scalar projection to the ideal preserves auxiliary degree, and
$\widetilde U_{n+1}$ intertwines its restriction in auxiliary degree
one with projection onto $M_{n+1}$.  The sum of the discarded squared
norms over the $z$ coordinates is therefore
\[
 \frac{n+1}{n+2}\langle A_n^Mq,q\rangle.
\]
Every product $y_bU_nq$ has auxiliary degree two and belongs to $J_M$,
so it is discarded in full.  On a normalized monomial of total degree
$\ell$ and auxiliary multi-degree $\beta$, the squared forward weight
in $y_b$ is $(\beta_b+1)/(\ell+1)$.  Summing over $b$ gives
$(|\beta|+r)/(\ell+1)$.  For each fixed coordinate, distinct monomials
have orthogonal images.  Hence the same identity holds on any vector
in a sector of fixed auxiliary degree.  On $E_n$ this sum is
$(r+1)/(n+2)$ times its squared norm.  Adding and polarizing proves
\eqref{eq:complete-defect}.

For a pure $z$ input, multiplication by an original coordinate remains
in the pure sector and has zero defect.  The sum of squared unprojected
auxiliary products is $r/(\ell+1)$ times the input norm; projection can
only decrease it.  For each coordinate, the discarded images of the
two input sectors have distinct auxiliary degrees, so their cross
inner products vanish.  This proves the reducing statement as well.
\end{proof}

\begin{proof}[Proof of Theorem~\ref{thm:main}]
Apply \eqref{eq:complete-scalar} to $J_M$ in $m+r$ variables, and
restrict it to $E_n$ in total degree $n+1$.  Dropping the positive
scalar summand in \eqref{eq:complete-defect} gives
\[
 \frac{n+1}{n+2}A_n^M\le\frac{C_{J_M}}{n+2}I,
 \qquad
 A_n^M\le\frac{C_{J_M}}{n+1}I.
\]
This is \eqref{eq:main-degree} in the original vector norm.
Lemma~\ref{op:commutators} proves
\eqref{eq:main-commutator} and the restriction block bound.
Lemma~\ref{op:hilbert} and the singular-value calculation in
Lemma~\ref{op:counting} give all the claimed weak and strong Schatten
classes, including the defect column and the projection commutators.
The finite-dimensional quotient case was accounted for there.
Lemma~\ref{op:graded} identifies the closed graded formulation with
the polynomial formulation.
\end{proof}

The order of these steps is essential for the dimension-sensitive
conclusion.  The scalar algebra in Lemma~\ref{complete:ring} has
dimension $m$ even when $T$ has smaller support.  Counting in that
scalar algebra would discard the sharper exponent.  Equation
\eqref{eq:complete-defect} first returns the degree estimate to
$Q_n^M$, whose dimension is the Hilbert function of $T$.

The square-zero part also explains the relevance of nonreduced ideals
in the scalar theorem.  Replacing $J_M$ by its radical removes the
entire sector that encodes the vector quotient.  Algebraic support
alone does not retain the information needed for this transfer.

% END sections/conclusion.tex


% BEGIN sections/extensions.tex
\section{Other ball norms, analytic ideals, and the endpoint}
\label{sec:extensions}

\subsection{The standard kernel family}
For $\lambda>0$, let $H^{(\lambda)}$ be the Hilbert function space on
$\mathbb B_m$ with reproducing kernel
\[
 K_\lambda(z,w)=(1-\langle z,w\rangle)^{-\lambda}.
\]
Its orthogonal monomial norms are
\begin{equation}\label{eq:ext-norm}
 \|z^\alpha\|_\lambda^2
 =\frac{\alpha!\,\Gamma(\lambda)}{\Gamma(|\alpha|+\lambda)}.
\end{equation}
These formulas follow by the binomial expansion of the kernel.  The
cases $\lambda=1$, $m$, and $m+1$ are respectively the
Drury--Arveson, sphere Hardy, and normalized ball Bergman spaces.

\begin{corollary}\label{ext:kernel-family}
The quotient and restriction conclusions of Theorem~\ref{thm:main}
hold for the closure of the same homogeneous submodule in
$H^{(\lambda)}\otimes\C^r$, for every $\lambda>0$, with the same
dimension exponents.  The constants may depend on $\lambda$.
\end{corollary}
\begin{proof}
The inner products on each homogeneous degree differ by a scalar,
so the degreewise unitary between normalized monomial bases maps the
ideal or submodule spaces, and their orthogonal complements, onto
one another.  Under this identification the quotient coordinate in
degree $n\to n+1$ is
\[
 S_{j,n}^{(\lambda)}=r_nS_{j,n},\qquad
 r_n=\sqrt{\frac{n+1}{n+\lambda}}.
\]
For $n\ge1$ this gives
\begin{align*}
 [(S_i^{(\lambda)})^*,S_j^{(\lambda)}]|_{Q_n}
 &=r_n^2[S_i^*,S_j]|_{Q_n}
   +(r_n^2-r_{n-1}^2)S_jS_i^*|_{Q_n},\\
 r_n^2-r_{n-1}^2
 &=\frac{\lambda-1}{(n+\lambda)(n+\lambda-1)}=O(n^{-2}).
\end{align*}
The sequence $r_n$ is bounded and the quotient coordinates in the
Drury--Arveson norm are contractions.  Hence the new commutator blocks
have norm $O((n+1)^{-1})$.  The same calculation applies to the
restriction coordinates.  The defect blocks are multiplied by $r_n$,
so their norms remain $O((n+1)^{-1/2})$.  The finitely many initial
terms are harmless, and the original Hilbert function supplies the
same singular-value counts.  Finite quotients remain finite.
\end{proof}

This transfer uses actual scalar weights between adjacent degrees.
It does not appeal to invariance of essential normality under an
arbitrary equivalent renorming.

\subsection{Ideals with holomorphic coefficients}
Ordinary homogeneity is particularly restrictive for a holomorphic
function at the origin.

\begin{proposition}\label{ext:holomorphic}
If a nonzero holomorphic function near $0\in\C^m$ satisfies
$f(tz)=t^df(z)$ for positive $t$ near $1$, then $d$ is a nonnegative
integer and $f$ is a homogeneous polynomial of degree $d$.
If $p_1,\ldots,p_s$ are homogeneous polynomials,
$J=(p_1,\ldots,p_s)\subset R$, and
$\mathfrak a=\sum_jp_j\cO(\mathbb B_m)$, then
\[
 \overline{\mathfrak a\cap H^{(\lambda)}}^{\,H^{(\lambda)}}
   =\overline J^{\,H^{(\lambda)}}.
\]
\end{proposition}
\begin{proof}
Differentiating in $t$ at $1$ gives
$\sum_jz_j\partial_jf=df$.  For the convergent Taylor expansion
$f=\sum_\alpha c_\alpha z^\alpha$, uniqueness of coefficients gives
$(|\alpha|-d)c_\alpha=0$.  Some coefficient is nonzero, so
$d\in\N$, and only the finitely many monomials of that degree occur.
Equality extends throughout a connected domain of holomorphy by
the identity theorem.

For the second assertion, write $F=\sum_jp_jg_j\in H^{(\lambda)}$
with the $g_j$ holomorphic.  If $\deg p_j=d_j$, the homogeneous
component is $F_n=\sum_jp_j(g_j)_{n-d_j}\in J_n$.  Its finite
homogeneous sums belong to $J$ and converge in the Hilbert norm,
because the homogeneous spaces are orthogonal.  Thus
$\mathfrak a\cap H^{(\lambda)}\subset\overline J$.
The reverse inclusion after closure follows from
$J\subset\mathfrak a\cap H^{(\lambda)}$.
\end{proof}

The coefficients $g_j$ in this proposition need not be bounded
multipliers.  If one starts with infinitely many homogeneous
holomorphic generators, the polynomial ring is Noetherian, so a
finite subset generates the same polynomial ideal and its analytic
extension.  The same component argument applies to homogeneous
vector-valued generators and to global sections of the generated
analytic ideal sheaf, since a representation near the origin already
determines every Taylor component.

\subsection{Why the endpoint is stated weakly}
The weak exponent cannot in general be replaced by its strong
counterpart, even for a simple homogeneous ideal.

\begin{example}\label{ext:sharp}
Let $I=(x^2)\subset\C[x,y]$.  Its quotient has support dimension one.
In degree $n\ge1$ an orthonormal basis is
\[
 e_n=y^n,\qquad f_n=\sqrt n\,xy^{n-1},
\]
and degree zero is spanned by $e_0=1$.  The compressed $x$ coordinate
satisfies
\[
 S_xe_n=\frac1{\sqrt{n+1}}f_{n+1},\qquad S_xf_n=0,
\]
and therefore
\[
 [S_x^*,S_x]e_n=\frac1{n+1}e_n,
 \qquad [S_x^*,S_x]f_n=-\frac1n f_n\quad(n\ge1).
\]
The eigenvalue on $e_0$ is $1$.  Its singular values thus belong to
$\ell^{1,\infty}$ but not to $\ell^1$.  In particular the quotient
commutator is in $\cS_{1,\infty}$ and in every $\cS_p$, $p>1$,
but is not trace class.  This example retains a transverse first
order nilpotent direction; replacing the ideal by its radical would
remove that direction.
\end{example}

\subsection{Exact defect blocks for homogeneous ideals}

The following calculations test the degree estimate and its normalization
without using the projective lifting theorem. They also exhibit the
necessity of the order $(n+1)^{-1}$ in the general defect-square bound.

\begin{example}[An ideal with an embedded component]\label{ext:embedded-defect}
Let $I=(x^2,xy)\subset\C[x,y,z]$. Its radical is $(x)$, while the class
of $x$ in $R/I$ is nonzero and has annihilator $(x,y)$. In particular
$(x,y)$ is an embedded associated prime. For $n\geq1$, an orthogonal
monomial basis of $Q_n$ consists of
\[
 z^n,\quad y^a z^{n-a}\ (1\leq a\leq n),\quad xz^{n-1}.
\]
Normalize these monomials in the Drury--Arveson norm. In any coordinate
$z_j$, the squared forward weight of a normalized degree-$n$ monomial
$z^\alpha$ is $(\alpha_j+1)/(n+1)$. Distinct monomials have orthogonal
images under each fixed coordinate, and projection onto the ideal
keeps exactly the monomials in the ideal. Thus each $B_{j,n}^*B_{j,n}$
is diagonal in this basis. No product of $z^n$ by a coordinate lies in
$I$. For $y^a z^{n-a}$, $a\geq1$, only the product by $x$ lies in $I$;
its discarded squared weight is $1/(n+1)$. For $xz^{n-1}$, the products
by $x$ and $y$ lie in $I$, with discarded squared weights $2/(n+1)$
and $1/(n+1)$, respectively. Consequently the eigenvalues of $A_n$,
with multiplicities indicated, are
\begin{equation}\label{eq:embedded-spectrum}
 \begin{array}{c|ccc}
 \text{eigenvalue}&0&1/(n+1)&3/(n+1)\\
 \text{multiplicity}&1&n&1.
 \end{array}
\end{equation}
In particular $\|A_n\|=3/(n+1)$ and $\Tr A_n=(n+3)/(n+1)$.
The trace agrees with \eqref{eq:op-trace}, since $h_n=n+2$ for
$n\geq1$. At degree zero $A_0=0$, as all coordinates survive.
This example has $\delta=2$ and shows that one cannot replace
$O((n+1)^{-1})$ by $o((n+1)^{-1})$ for every homogeneous ideal.
\end{example}

\begin{example}[A nonmonomial quadric]\label{ext:quadric-defect}
Let $m\geq3$, $q(z)=\sum_{j=1}^m z_j^2$, and $I=(q)$. On the
polynomial ring equip monomials with the Fischer inner product
$\langle z^\alpha,z^\beta\rangle_{\rm F}=\alpha!\,\delta_{\alpha\beta}$.
Within degree $n$ it is $n!$ times the Drury--Arveson inner product;
therefore the orthogonal projections coincide. The Fischer adjoint of
multiplication by $q$ is $\Delta=\sum_j\partial_{z_j}^2$, so
$Q_n=\ker(\Delta|_{R_n})$. Take $p\in Q_n$, $n\geq1$. The projection
of $z_jp$ onto $I_{n+1}=qR_{n-1}$ is
\begin{equation}\label{eq:quadric-projection}
 B_{j,n}p=\frac{q\,\partial_{z_j}p}{2n+m-2}.
\end{equation}
Indeed, $\partial_{z_j}p$ is harmonic and homogeneous of degree $n-1$.
The product rule and Euler identity give
\[
 \Delta(z_jp)=2\partial_{z_j}p,\qquad
 \Delta(q\,\partial_{z_j}p)=2(2n+m-2)\partial_{z_j}p.
\]
Hence the right side of \eqref{eq:quadric-projection} belongs to
$I_{n+1}$ and its difference from $z_jp$ belongs to $Q_{n+1}$.
Fischer adjointness gives
\[
 \|q\,\partial_{z_j}p\|_{\rm F}^2
   =2(2n+m-2)\|\partial_{z_j}p\|_{\rm F}^2,
 \qquad
 \sum_j\|\partial_{z_j}p\|_{\rm F}^2=n\|p\|_{\rm F}^2,
\]
where the second equality is the Euler identity paired with $p$.
Converting the output squared norm by $(n+1)!$ and the input squared
norm by $n!$, we obtain
\[
 \sum_j\|B_{j,n}p\|_{\DA}^2
   =\frac{2n}{(n+1)(2n+m-2)}\|p\|_{\DA}^2.
\]
Polarization yields the exact positive block
\begin{equation}\label{eq:quadric-defect}
 A_n=\frac{2n}{(n+1)(2n+m-2)}I_{Q_n}\qquad(n\geq1).
\end{equation}
At $n=0$ all degree-one polynomials survive and $A_0=0$. In particular,
$n\|A_n\|\to1$. This gives a nonmonomial test of the adjacent-degree
conversion and again shows sharpness of the degree-decay order.

There is also an independent trace check of the coefficient. Since
multiplication by the nonzero quadratic polynomial is injective,
\[
 h_n=\binom{n+m-1}{m-1}-\binom{n+m-3}{m-1}
 =\frac{(2n+m-2)(n+m-3)!}{n!(m-2)!},
\]
where the second binomial is zero if $n<2$. Thus for $n\geq0$,
\[
 \frac{\Tr A_n}{h_n}
 =\frac{n+m}{n+1}-\frac{h_{n+1}}{h_n}
 =\frac{2n}{(n+1)(2n+m-2)}.
\]
The Fischer calculation established scalarity, so the trace gives
exactly the coefficient in \eqref{eq:quadric-defect}.
\end{example}

\begin{example}[The $2\times3$ determinantal ideal]\label{ext:segre-defect}
Let $R=\C[z_{ij}:1\leq i\leq2,\ 1\leq j\leq3]$ and let $I$ be the
ideal generated by all $2\times2$ minors of the matrix $(z_{ij})$.
The substitution $z_{ij}\mapsto u_i v_j$ induces in each degree $n$
an isomorphism
\[
 (R/I)_n\simeq \C[u_1,u_2]_n\otimes\C[v_1,v_2,v_3]_n.
\]
Indeed, each monomial of bidegree $(n,n)$ in the two sets of variables
is the image of a product of $n$ pairs $z_{ij}$. Two such products have
the same image precisely when the multisets of row and column indices
agree. After ordering the row indices, their column assignments are
related by transpositions. A transposition of column indices between
two different rows and two different columns is precisely the relation
$z_{ij}z_{k\ell}=z_{i\ell}z_{kj}$ modulo a minor; transpositions
with equal row or column indices leave the product unchanged. Hence
all products with the same image agree modulo $I$. Distinct image
monomials are linearly independent, proving the isomorphism and
\[
 h_n=(n+1)\binom{n+2}{2}=\frac{(n+1)^2(n+2)}2.
\]
The tensor product action of $U(2)\times U(3)$ on the six coordinate
variables is unitary. Its induced action on the Drury--Arveson space
is unitary, since unitary changes of variables preserve the kernel
$(1-\langle z,w\rangle)^{-1}$. The substitution is equivariant
for the natural actions on the two factors, so $I$ is invariant;
thus the quotient and its orthogonal realization $Q_n$ are equivalent
representations. This bidegree space is irreducible: an invariant
subspace is a sum of weight spaces for the diagonal tori, which are
one-dimensional monomial spaces. From any one monomial, the operators
$u_i\partial_{u_k}$ and $v_j\partial_{v_\ell}$ in the complexified
Lie algebra generate every monomial of the same bidegree by successive
transfers of exponents. Consequently any nonzero invariant subspace
is the whole bidegree space.

The operator $A_n$ is invariant under this action. Indeed the
orthogonal projections onto $I$ and $Q$ commute with it, and a unitary
linear change of the six coordinates rotates the column of discarded
coordinate maps by a unitary matrix on the column index. The sum of
their adjoint squares is unchanged. Schur's lemma now shows that
$A_n$ is scalar. Since $m=6$, the trace identity
\eqref{eq:op-trace} gives for every $n\geq0$
\begin{equation}\label{eq:segre-defect}
 A_n=\left(\frac{n+6}{n+1}
       -\frac{(n+2)(n+3)}{(n+1)^2}\right)I_{Q_n}
     =\frac{2n}{(n+1)^2}I_{Q_n}.
\end{equation}
Here $A_0=0$ and $n\|A_n\|\to2$. This calculation tests the rate for
a nonprincipal, nonmonomial ideal with nontrivial syzygies; it uses
no projective lifting estimate.
\end{example}

The theorem imposes no smoothness or transversality assumption on
the projective support or its link in the sphere.  Its analytic
estimates are consequences of the homogeneous presentation, not
additional hypotheses on that support.  The use of the standard
grading is nevertheless essential to the argument.  A
nonhomogeneous relation such as $z_2-z_1^2$ need not contain its
homogeneous components in the ideal.  Such an ideal does not give
the projective ideal sheaf with the exact degree decomposition used
in \eqref{eq:geo-sections}; the theorem above makes no conclusion
about that general situation.

% END sections/extensions.tex

\appendix

% BEGIN sections/classical_inputs.tex
\section{Classical inputs and their uses}
\label{sec:classical-inputs}

We record the imported results with their hypotheses.  None is an
essential-normality or stable-division theorem for the ideal under
consideration.

\subsection{Noetherianity and the Hilbert function}
The polynomial ring over a field is Noetherian, and every submodule of
a finite free module over a Noetherian ring is finitely generated
\cite{StacksNoetherian}.  This is used to recover polynomial relations
from a closed graded submodule, to choose finite presentations, and
to choose finite module Gr\"obner bases.

For a finite graded module $T$ over a standard polynomial ring, its
Hilbert function is eventually polynomial; for positive support
dimension $\delta$ the degree of that polynomial is $\delta-1$.
The module version follows from the graded prime filtration: its
successive quotients are $(R/\mathfrak p)(-a)$ with homogeneous
primes $\mathfrak p$.  Such a filtration is finite by Noetherianity.
Hilbert functions add in short exact sequences.  The largest dimension
of these prime quotients is $\dim\Supp T$, and the positive leading
coefficients of the corresponding Hilbert polynomials cannot cancel.
The ring case is the Hilbert--Serre dimension theorem
\cite{StacksHilbert}; the coherent-sheaf formulation is in
\cite{SerreFAC}. This input is used both in the preliminary trace
asymptotic \eqref{eq:op-trace-asymptotic} and, after the degree estimate
in the original quotient has been obtained, in the final Schatten
counting argument. The trace asymptotic is not an input to the proof of
the general degreewise operator-norm estimate.

\subsection{Regularity, Quot, and base change}
A coherent sheaf $\mathcal F$ on $\Pj^N$ is $D$-regular if
$H^j(\Pj^N,\mathcal F(D-j))=0$ for all $j>0$.
Regularity persists in larger degrees, implies global generation at
degree $D$, and makes multiplication by linear forms surjective on
all subsequent section spaces.  For fixed $N,a$, and a Hilbert
polynomial $P$, the kernels of quotients
$\cO^{\oplus a}\twoheadrightarrow\mathcal F$ with polynomial $P$
have a common regularity bound.  A common line twist of the source
changes only the degree index.  The precise results are
\cite[Tags 08A6--08A8 and 08AG]{StacksRegularity}.

The quotients of a fixed coherent sheaf on projective space with
fixed Hilbert polynomial are parametrized by a projective Quot
scheme, carrying a universal flat quotient
\cite[Theorem~3.2]{GrothendieckHilbert}; properness in the case used
here is also \cite[Tag 0DPA]{StacksQuot}.  The kernel of the universal
quotient is flat over the base, because the middle term and quotient
are flat.  The degree of evaluation is chosen uniformly over this
whole family.

For a proper morphism to a locally Noetherian base and a coherent
sheaf flat over that base, vanishing of the first cohomology on each
fiber implies that the direct image of the sheaf is locally free
and that its degree-zero cohomology commutes with base change
\cite[Theorem~1.1 and Corollary~1.2]{ConradBaseChange}.
No reducedness assumption on the parameter scheme is required.
Applied after the chosen regular twist, this identifies the
complete section spaces as fibers of a vector bundle and supplies
the Grassmannian morphism used in the compactness argument.

\subsection{Analytification and Stein exactness}
On complex projective space, analytification is exact for coherent
algebraic sheaves and identifies their algebraic and analytic global
sections.  Conversely, coherent analytic sheaves on projective space
are algebraizable.  These are the GAGA comparisons
\cite[Proposition~10 and Theorems~1--3]{SerreGAGA}.
They allow one exact coherent presentation to be used both in the
algebraic limit argument and in holomorphic division on coordinate
balls.

For a coherent analytic sheaf $\mathcal F$ on a Stein space $U$,
Cartan's theorem B asserts $H^j(U,\mathcal F)=0$ for $j>0$
\cite{GrauertRemmert}.  In the division theorem the exact sequence
\[
 0\longrightarrow\ker A\longrightarrow\cO_U^a
   \longrightarrow\im A\longrightarrow0
\]
therefore makes the map on global sections onto its image sheaf.
It supplies an initial holomorphic lift on a ball.  Its norm is not
estimated and is not used in the final constant: polynomial
approximations of this lift are divided again using the uniform
coefficient estimate.  Cartan's theorem is consequently a
qualitative input at that step.

\subsection{The positive-curvature solution theorem}
Let $(W,\omega)$ be a complete K\"ahler manifold and $F$ a Hermitian
holomorphic vector bundle.  If the curvature operator
$[i\Theta(F),\Lambda_\omega]$ on $(p,q)$-forms, $q\ge1$, is bounded
below by $\Delta I$ with $\Delta>0$, then a closed $L^2$ form $g$
of that type admits a solution $\bar\partial u=g$ with
\[
 \|u\|^2\le\Delta^{-1}\|g\|^2.
\]
This is the standard Bochner--Kodaira--Nakano $L^2$ solution theorem
\cite[Theorem~6.1]{DemaillyL2}; see also \cite{Demailly}.  We apply it on compact $W=\Pj^N$, to
$(N,q)$-forms with coefficient bundle
$E_\ell\otimes L^h\otimes K_W^{-1}$.  The identification with
$(0,q)$-forms valued in $E_\ell\otimes L^h$ preserves the
Dolbeault operator and the induced norms.  The proof in the text
establishes the required lower bound
$\Delta\ge c_Lh-C_0-C_1\varepsilon^{-2}$ and chooses the scale
so that $\Delta$ is a positive multiple of $n$.

Smooth data give smooth minimal solutions by elliptic regularity
for the Dolbeault Laplacian on a compact manifold.  The quantitative
bound uses only the positive curvature gap; no uniform higher
regularity estimate is needed in the finite descent.

The remaining estimates---the preservation of image dimensions at
normalized limits, uniform two-ball division, the relative metric
derivative bounds, overlap comparisons, and the finite correction
with its constants---are proved in the body of the article.

\subsection{Proof dependencies and scope}
For clarity, the quantitative proof uses the following implications. The
arrows below indicate the precise outputs consumed by the next stage.
\begin{enumerate}
\item Lemma~\ref{geo:norm} and saturation identify the original degree
projections with projective projections. They do not provide any division
bound.
\item Theorem~\ref{thm:B} retains the image Hilbert functions of the
normalized matrices. Its proof uses complete evaluation in a degree
chosen for the entire proper parameter family, not convergence of the
normalizing changes of basis. Sheaf exactness is required only at
$E_0,\ldots,E_{N-1}$, together with generation of the first-row image
ideal sheaf.
\item Theorem~\ref{div:uniform} gives uniform Euclidean division on nested
balls for those matrices. The datum is required to be in the analytic
image sheaf on the whole input ball. The stronger all-degree Hilbert
function hypothesis supplied by B suffices; the proof also establishes
the weaker eventual-polynomial variant.
\item Theorem~\ref{thm:D} bounds curvature relative to the same metrics
that normalize division. Comparability to sampled metrics holds on
the smaller buffered charts $V_\alpha$, which contain all the local
balls once the scale is sufficiently small. The microscopic estimate
\eqref{ana:micro-bound} converts C to those metrics, including the original
line norm. The curvature estimate and choice of scale give the
$O(\varepsilon)$ Dolbeault solution bound \eqref{ana:solver}.
\item The buffered cochain construction uses holomorphic division only
on holomorphic data in degrees through $N$; the finite descent uses
$\bar\partial$ only on closed data. In particular the top-degree form
is closed by dimension, not by injectivity of $d_N$.
Proposition~\ref{ana:eta-identities} and estimate
\eqref{ana:final-lift-bound} give the exact equation and degree-independent
bound of Theorem~\ref{geo:lifting}.
\item Proposition~\ref{geo:transfer} consumes the derivative estimate of D
and the lift of E in the same coefficient metric. The original
Drury--Arveson degree normalization removes the auxiliary metric.
The identity \eqref{eq:complete-defect} then transfers the bound to the
original vector quotient. Hilbert--Serre counting proves A.
\end{enumerate}
The analytic part makes no assertion of a bounded division operator in
fixed coefficient metrics, nor of an intrinsic restriction measure on a
singular variety. No smoothness, normality, radicality, or boundary
transversality hypothesis on the scheme has been used. The proof concerns
standard homogeneous relations in finitely many variables and finite
coefficient multiplicity. It does not establish the general
nonhomogeneous or unequal-weight graded conjectures. The manuscript's
proof and historical positioning are offered for independent mathematical
assessment.

% END sections/classical_inputs.tex

\raggedbottom
\begin{thebibliography}{99}

\bibitem{AnderssonMembership}
M.~Andersson, The membership problem for polynomial ideals in terms of
residue currents, \emph{Ann. Inst. Fourier} \textbf{56} (2006), no.~1,
101--119. \url{https://www.numdam.org/item/AIF_2006__56_1_101_0/}.

\bibitem{AnderssonWulcan}
M.~Andersson and E.~Wulcan, Residue currents with prescribed annihilator
ideals, \emph{Ann. Sci. \'Ec. Norm. Sup. (4)} \textbf{40} (2007), no.~6,
985--1007. \url{https://arxiv.org/abs/math/0511237}.

\bibitem{Arveson1998}
W.~Arveson, Subalgebras of $C^*$-algebras III: multivariable operator
theory, \emph{Acta Math.} \textbf{181} (1998), no.~2, 159--228.
\url{https://arxiv.org/abs/funct-an/9705007}.

\bibitem{Arveson2005}
W.~Arveson, $p$-summable commutators in dimension $d$,
\emph{J. Operator Theory} \textbf{54} (2005), no.~1, 101--117.
\url{https://arxiv.org/abs/math/0308104}.

\bibitem{Beatrous}
F.~Beatrous, Jr., $L^p$-estimates for extensions of holomorphic functions,
\emph{Michigan Math. J.} \textbf{32} (1985), 361--380.
\href{https://doi.org/10.1307/mmj/1029003244}{doi:10.1307/mmj/1029003244}.

\bibitem{ConradBaseChange}
B.~Conrad, \emph{Applications of base change for coherent cohomology},
Math 248B notes, Theorem~1.1 and Corollary~1.2.
\url{https://math.stanford.edu/~conrad/248BPage/handouts/cohom.pdf}.

\bibitem{Demailly}
J.-P.~Demailly, \emph{Complex Analytic and Differential Geometry},
online monograph, version of 21 June 2012, Chapters VI--VIII.
\url{https://www-fourier.univ-grenoble-alpes.fr/~demailly/manuscripts/agbook.pdf}.

\bibitem{DemaillyL2}
J.-P.~Demailly, \emph{Applications of the theory of $L^2$ estimates and
positive currents in algebraic geometry}, lecture notes, 2007,
Section~6, Theorem~6.1.
\url{https://www-fourier.univ-grenoble-alpes.fr/~demailly/manuscripts/eem2007.pdf}.

\bibitem{DouglasIndex}
R.~G.~Douglas, A new kind of index theorem, in
\emph{Analysis, Geometry and Topology of Elliptic Operators},
World Scientific, Hackensack, NJ, 2006, pp.~369--382.
\url{https://arxiv.org/abs/math/0507542}.

\bibitem{DouglasJabbariTangYu}
R.~G.~Douglas, M.~Jabbari, X.~Tang, and G.~Yu, A new index theorem for
monomial ideals by resolutions, \emph{J. Funct. Anal.} \textbf{275}
(2018), 735--760. \url{https://arxiv.org/abs/1708.05796}.

\bibitem{DouglasTangYu}
R.~G.~Douglas, X.~Tang, and G.~Yu, An analytic Grothendieck Riemann Roch
theorem, \emph{Adv. Math.} \textbf{294} (2016), 307--331.
\url{https://arxiv.org/abs/1404.4396}.

\bibitem{EnglisEschmeier}
M.~Engli\v{s} and J.~Eschmeier, Geometric Arveson--Douglas conjecture,
\emph{Adv. Math.} \textbf{274} (2015), 606--630;
corrigendum, \emph{Adv. Math.} \textbf{278} (2015), 254.
\href{https://doi.org/10.1016/j.aim.2014.11.026}{doi:10.1016/j.aim.2014.11.026}.
Authors' preprint:
\url{https://users.math.cas.cz/~englis/85.pdf}.

\bibitem{EschmeierCriterion}
J.~Eschmeier, Essential normality of homogeneous submodules,
\emph{Integral Equations Operator Theory} \textbf{69} (2011), no.~2,
171--182.
\href{https://doi.org/10.1007/s00020-010-1809-4}{doi:10.1007/s00020-010-1809-4}.

\bibitem{Galligo}
A.~Galligo, Th\'eor\`eme de division et stabilit\'e en g\'eom\'etrie
analytique locale, \emph{Ann. Inst. Fourier} \textbf{29} (1979), no.~2,
107--184. \href{https://doi.org/10.5802/aif.745}{doi:10.5802/aif.745}.

\bibitem{GrauertRemmert}
H.~Grauert and R.~Remmert, \emph{Theory of Stein Spaces},
Grundlehren der mathematischen Wissenschaften, vol.~236,
Springer, New York, 1979.
\href{https://doi.org/10.1007/978-1-4757-4357-9}{doi:10.1007/978-1-4757-4357-9}.

\bibitem{GrothendieckHilbert}
A.~Grothendieck, Techniques de construction et th\'eor\`emes
d'existence en g\'eom\'etrie alg\'ebrique IV: les sch\'emas de Hilbert,
\emph{S\'eminaire Bourbaki}, expos\'e 221 (1961), 249--276.
\url{https://numdam.org/item/SB_1960-1961__6__249_0/}.

\bibitem{GuoWang}
K.~Guo and K.~Wang, Essentially normal Hilbert modules and
$K$-homology, \emph{Math. Ann.} \textbf{340} (2008), no.~4,
907--934.
\href{https://doi.org/10.1007/s00208-007-0175-2}{doi:10.1007/s00208-007-0175-2}.

\bibitem{Hartshorne}
R.~Hartshorne, \emph{Algebraic Geometry}, Graduate Texts in
Mathematics, vol.~52, Springer, New York, 1977.

\bibitem{HauserMueller}
H.~Hauser and G.~M\"uller, A rank theorem for analytic maps between power
series spaces, \emph{Publ. Math. IH\'ES} \textbf{80} (1994), 95--115.
\url{https://www.numdam.org/item/PMIHES_1994__80__95_0/}.

\bibitem{HormanderGenerators}
L.~H\"ormander, Generators for some rings of analytic functions,
\emph{Bull. Amer. Math. Soc.} \textbf{73} (1967), 943--949.
\href{https://doi.org/10.1090/S0002-9904-1967-11860-3}{doi:10.1090/S0002-9904-1967-11860-3}.

\bibitem{Ji}
Q.~Ji, Division theorems for exact sequences, \emph{Math. Ann.}
\textbf{355} (2013), 947--984. \url{https://arxiv.org/abs/1102.3950}.

\bibitem{KelleherTaylor}
J.~J.~Kelleher and B.~A.~Taylor, Finitely generated ideals in rings of
analytic functions, \emph{Math. Ann.} \textbf{193} (1971), 225--237.

\bibitem{SerreFAC}
J.-P.~Serre, Faisceaux alg\'ebriques coh\'erents,
\emph{Ann. of Math. (2)} \textbf{61} (1955), 197--278.

\bibitem{SerreGAGA}
J.-P.~Serre, G\'eom\'etrie alg\'ebrique et g\'eom\'etrie analytique,
\emph{Ann. Inst. Fourier} \textbf{6} (1956), 1--42.
\url{https://numdam.org/item/AIF_1956__6__1_0/}.

\bibitem{ShalitDivision}
O.~M.~Shalit, Stable polynomial division and essential normality of
graded Hilbert modules, \emph{J. Lond. Math. Soc. (2)}
\textbf{83} (2011), no.~2, 273--289.
\href{https://doi.org/10.1112/jlms/jdq054}{doi:10.1112/jlms/jdq054}.
\url{https://arxiv.org/abs/1003.0502}.

\bibitem{SkodaDivision}
H.~Skoda, Application des techniques $L^2$ \`a la th\'eorie des id\'eaux
d'une alg\`ebre de fonctions holomorphes avec poids,
\emph{Ann. Sci. \'Ec. Norm. Sup. (4)} \textbf{5} (1972), 545--579.
\url{https://www.numdam.org/item/ASENS_1972_4_5_4_545_0/}.

\bibitem{SkodaBundles}
H.~Skoda, Morphismes surjectifs de fibr\'es vectoriels semi-positifs,
\emph{Ann. Sci. \'Ec. Norm. Sup. (4)} \textbf{11} (1978), 577--611.
\href{https://doi.org/10.24033/asens.1357}{doi:10.24033/asens.1357}.

\bibitem{StacksHilbert}
The Stacks Project Authors, \emph{The Stacks Project},
Sections 10.58 and 10.117, especially Lemma 10.117.1 (Tag 00P6),
and Section 33.45, numerical intersections and dimensions of support.
\url{https://stacks.math.columbia.edu/tag/00P6}.

\bibitem{StacksNoetherian}
The Stacks Project Authors, \emph{The Stacks Project},
Section 10.31, Lemmas 10.31.1 and 10.31.4, Noetherian rings and modules.
\url{https://stacks.math.columbia.edu/tag/00FM}.

\bibitem{StacksQuot}
The Stacks Project Authors, \emph{The Stacks Project},
Lemma 108.6.1 (Tag 0DPA), properness of the fixed-polynomial Quot space.
\url{https://stacks.math.columbia.edu/tag/0DPA}.

\bibitem{StacksRegularity}
The Stacks Project Authors, \emph{The Stacks Project},
Lemmas 33.35.10--12 (Tags 08A6--08A8) and 33.35.18 (Tag 08AG),
Castelnuovo--Mumford regularity and its uniform bound for quotient kernels.
\url{https://stacks.math.columbia.edu/tag/08AG}.

\bibitem{Varolin}
D.~Varolin, Division theorems and twisted complexes,
\emph{Math. Z.} \textbf{259} (2008), 1--20.
\href{https://doi.org/10.1007/s00209-007-0133-4}{doi:10.1007/s00209-007-0133-4}.

\bibitem{YiWang}
Y.~Wang, Essential normality---a unified approach in terms of local
decompositions, \emph{Proc. Lond. Math. Soc. (3)} \textbf{119} (2019),
1388--1429. \url{https://arxiv.org/abs/1812.03211}.

\end{thebibliography}
\end{document}
